%% Apuntes de Optica Cu\'antica
%
%\documentclass[twocolumn]{article}
\documentclass[twocolumn,aps,eqsecnum]{revtex4}
\usepackage[spanish,activeacute]{babel}
\usepackage{amsmath} 
\usepackage{cancel}
\usepackage{color}
 
% vectors (the most common ones)
\newcommand{\mbd}{\mathbf{d}}
\newcommand{\mbe}{\mathbf{e}}
\newcommand{\mbi}{\mathbf{i}}
\newcommand{\mbj}{\mathbf{j}}
\newcommand{\mbk}{\mathbf{k}}
\newcommand{\mbp}{\mathbf{p}} 
\newcommand{\mbq}{\mathbf{q}}
\newcommand{\mbr}{\mathbf{r}}
%\newcommand{\mbrp}{\mathbf{r}^{\prime}}
\newcommand{\mbv}{\mathbf{v}}
\newcommand{\mbB}{\mathbf{B}}
\newcommand{\mbE}{\mathbf{E}}
\newcommand{\mbR}{\mathbf{R}}

\newcommand{\mcE}{\mathcal{E}} 
\newcommand{\mcH}{\mathcal{H}} 	%Hamiltonian 
\newcommand{\mcL}{\mathcal{L}} 
\newcommand{\mcV}{\mathcal{V}} 	%Hamiltonian
\newcommand{\Tr}{\mathrm{Tr}}
\newcommand{\qed}{\mathrm{Q.E.D.}}
\newcommand{\p}{\partial}
%\newcommand{\lan}{\langle}     \newcommand{\ran}{\rangle}
\newcommand{\adg}{a^{\dag}}
\newcommand{\Adg}{A^{\dag}} 	
\newcommand{\Hdg}{H^{\dag}} 	 
\newcommand{\Udg}{U^{\dag}} 	
%\newcommand{\lan}{\langle}
%\newcommand{\ran}{\rangle}


\providecommand{\abs}[1]{\left\lvert#1\right\rvert}
\providecommand{\norm}[1]{\lVert #1 \rVert}
\providecommand{\vxp}[1]{\langle #1 \rangle}
\providecommand{\bra}[1]{\langle #1 \rvert}
\providecommand{\ket}[1]{\lvert #1 \rangle}
\providecommand{\bbra}[1]{\langle #1 \rvert\rvert}
\providecommand{\kket}[1]{\lvert\lvert #1 \rangle}
\providecommand{\kkket}[1]{\lvert #1 \rangle\rangle}
\providecommand{\braket}[2]{\langle #1 \rvert #2 \rangle}
%\providecommand{\vxp3}[3]{\langle #1 |#2| #3 \rangle}
\providecommand{\tr}[1]{\text{tr}\left[ #1 \right]}
\newcommand{\ketbra}[2]{\left| {#1} \right\rangle\left\langle {#2}\right|}
%\newcommand{\mmax}{_{\text{max}}}

\newcommand{\mru}[1]{\mathrm{#1}} %unidades fisicas
\newcommand{\inc}[1]{\mathrm{(#1)}} %incisos en formulas desplegadas
\newcommand{\tbn}[1]{\textbf{#1}} %numero de problema

\begin{document}

\title{Apuntes de \'Optica Cu\'antica Semicl\'asica}
\author{H\'ector M. Castro Beltr\'an}
\affiliation{CIICAp-UAEM}
\date{\today}
\maketitle


Textos: Carmichael 1; Meystre \& Sargent III (MS); Scully \& Zubairy (SZ); 
Loudon; Stenholm; Berman \& Malinovsky (BM), Sargent III, Scully \& Lamb 
(SSL);  Fox. El n\'umero despu\'es de MS y Loudon identifica la edici\'on de 
 donde fue tomado el apunte. 

Los temas de Teor\'ia de fotodetecci\'on est\'an en el apunte 
``notas Fotodetecci\'on.pdf''


%%%%%%%%%%%%%%%%%%%%%%%%%%%%%%%%%%%
%%%%%%%%%%%%%%%%%%%%%%%%%%%%%%%%%%%
\section{Interacci\'on \'Atomo-Campo EM} 
En \'optica cu\'antica (OC) nos interesan procesos en donde la radiaci\'on, la 
materia o ambas deben describirse por la mec\'anica cu\'antica (MC). En estas 
notas consideramos la \textit{teor\'ia semicl\'asica}, en donde la luz es tratada 
por la teor\'ia de Maxwell y la materia mediante la MC, y la \textit{teor\'ia 
cu\'antica}, en donde luz y materia son ambas tratadas por la MC. No 
consideramos casos en los que la luz es cuantizada y la materia es tratada cl\'asicamente. 

Estudiaremos varios procesos, m\'as o menos en orden de complejidad 
creciente, tanto del sistema estudiado como por los m\'etodos. Por ejemplo, 
iniciaremos con sistemas donde no hay p\'erdidas de energ\'ia o coherencia, 
descrita por la funci\'on de onda, y luego pasaremos a sistemas disipativos 
y/o incoherentes, descritos por operador de densidad. Por simplicidad, el 
t\'ermino \'atomo incluye tambi\'en mol\'eculas, u otros entes materiales. 

En un curso introductorio de MC nos enfocamos en conocer los niveles de 
energ\'ia de un \'atomo. Esto no nos dice en realidad nada de los procesos 
por los cuales un \'atomo cambia de nivel (absorci\'on, emisi\'on). Para ello se 
requiere la interacci\'on con otro sistema, que puede ser luz u otros \'atomos 
(p. ej., colisiones). Tradicionalmente, estas interacciones se estudian 
primeramente si son d\'ebiles, donde la energ\'ia de interacci\'on es mucho 
mas peque\~na que la energ\'ia entre los niveles que se desea estudiar. Este 
es el dominio de la teor\'ia de perturbaciones. Consideraremos algunos de 
estos procesos, aunque no es mas dif\'icil plantear primero una situaci\'on mas 
general.  

\subsection{Efecto Fotoel\'ectrico} 
El Efecto Fotoel\'etrico (EF) es el proceso por el cual un \'atomo o material 
que absorbe luz de frecuencia mayor que la frecuencia de ionizaci\'on o de 
trabajo expulsa un electr\'on. En un fotodetector este electr\'on es sujeto a 
un proceso de amplificaci\'on para poder ser analizado por varios 
instrumentos electr\'onicos (osciloscopios, etc).  Aqu\'i describimos el proceso 
de excitaci\'on del \'atomo por un campo EM, iniciando con un repaso breve 
de algunos procesos de interacci\'on radiaci\'on materia hasta llegar al 
an\'alisis semicl\'asico del EF. En secciones posteriores estudiaremos otros 
aspectos de la teor\'ia de fotodetecci\'on. 

\subsection{Ecuaciones para las Amplitudes} 
Estudiamos la interacci\'on de un sistema no perturbado con Hamiltoniano 
$H_0$ con otro dado por el potencial $V$ que induce transiciones entre los 
estados del sistema no perturbado.  Presentaremos el c\'alculo en las 
im\'agenes de Schr\"odinger y de Interacci\'on. Posteriormente exploraremos 
estos conceptos. 

\textit{Imagen de Interacci\'on}: La funci\'on de onda es 
%
\begin{eqnarray} 
\Psi(\mbr,t) &=& \sum_n C_n (t) u_n (\mbr) e^{-i \omega_n t} \nonumber \\ 
&=& \ket{\Psi} = \sum_n C_n (t)  e^{-i \omega_n t} \ket{n} \,, 
\end{eqnarray}
%% 
donde $C_n (t)$ son llamadas \textit{amplitudes de probabilidad}, que llevan 
la informaci\'on temporal de la interacci\'on; $u_n (\mbr) = \ket{n}$ son las 
funciones de onda que obedecen la ecuaci\'on de Schr\"odinger independiente 
del tiempo $H_0 \ket{n} = \hbar \omega_n \ket{n}$; y $\hbar \omega_n$ es la 
energ\'ia del nivel $n$. La interacci\'on esta descrita por el operador $V$, que 
conecta dos o m\'as niveles del sistema. Tenemos entonces la ecuaci\'on de 
Schr\"odinger dependiente del tiempo 
$i \hbar \ket{\dot{\Psi}} =  (H_0 +V) \ket{\Psi}$ ($\dot{x} = dx/dt $), 
que nos llevar\'a a un conjunto de ecuaciones para las amplitudes que describen el acoplamiento en el tiempo entre dos niveles:   
%
\begin{eqnarray} 
\sum_n ( i \hbar \dot{C}_n + \cancel{\hbar \omega_n C_n} ) 
	e^{-i \omega_n t} \ket{n} 
	&=& \sum_n ( \cancel{\hbar \omega_n} +V) C_n 
	e^{-i \omega_n t} \ket{n} 	\nonumber \\ 
\sum_n i \hbar \dot{C}_n  e^{-i \omega_n t} \ket{n} 
	&=& \sum_n V C_n e^{-i \omega_n t} \ket{n} \,.
\end{eqnarray}
%% 
Multiplicando por $\bra{m}$ por la izquierda,  
%
\begin{eqnarray} 	\label{eq:amplitudeEqsIP}
\sum_n i \hbar \dot{C}_n  e^{-i \omega_n t}  \braket{m}{n} 
	&=& \sum_n C_n e^{-i \omega_n t}  \bra{m} V \ket{n} 	\nonumber \\ 
i \hbar \dot{C}_m  e^{-i \omega_m t}  
	&=& \sum_n C_n e^{-i \omega_n t} V_{mn} 	\nonumber \\ 
\dot{C}_m &=& - (i/\hbar) \sum_n C_n e^{i \omega_{mn} t} V_{mn} \,,
	\nonumber \\
\end{eqnarray}
%% 
donde $\omega_{mn} = \omega_m - \omega_n$, 
$\braket{m}{n} = \delta_{mn}$ y $ \bra{m} V \ket{n} = V_{mn}$. 

\textit{Imagen de Schr\"odinger}: La funci\'on de onda es 
%
\begin{eqnarray} 
\Psi(\mbr,t) &=& \sum_n c_n (t) u_n (\mbr) \,,	\qquad 
	 \ket{\Psi} = \sum_n c_n (t)  \ket{n} \,,		\nonumber \\ 
 c_n (t) &=& C_n e^{i \omega_n t} 	\,.
\end{eqnarray}
%% 
Substituyendo en la ecuaci\'on de onda 
%
\begin{eqnarray} 
i \hbar \sum_n  \dot{c}_n \ket{n} 
	&=& \sum_n (\hbar \omega_n +V) c_n  \ket{n} 	
\end{eqnarray}
%% 
y multiplicando por $\bra{m}$ por la izquierda, 
%
\begin{eqnarray} 	\label{eq:amplitudeEqsSP}
\sum_n i \hbar \dot{c}_n  \braket{m}{n} 
	&=&  \sum_n \hbar \omega_n c_n \braket{m}{n}
	 + \sum_n c_n  \bra{m} V \ket{n}  \nonumber \\ 
\dot{c}_m  &=& -i \omega_m c_m - (i/\hbar)  \sum_n c_n V_{mn} \,.
\end{eqnarray}
%% 
Hasta aqu\'i el c\'alculo es exacto, es decir, no hemos usado la teor\'ia 
perturbativa. 


%%%%%%%%%%%%%%%%%%%%%%%%%%%%%%%%%%%
\subsection{Perturbaciones Dependientes del Tiempo} 
De manera similar al caso de perturbaciones independientes del tiempo 
introducimos el par\'ametro $\lambda$, de modo que el Hamiltoniano se 
escribe como $H = H_0 + \lambda V(t)$, la expansi\'on en serie de los 
coeficientes como ($C_n \to a_n$) 
$a_n (t) = a_n^{(0)} + \lambda a_n^{(1)} + \lambda^2 a_n^{(2)} + \cdots$ 
y su substituci\'on en  
$\dot{a}_m = - (i/\hbar) \sum_n a_n e^{-i \omega_{mn} t} V_{mn}$ (imagen 
de Interacci\'on) nos lleva a  
%
\begin{eqnarray} 
\lefteqn{ \dot{a}_m^{(0)} + \lambda \dot{a}_m^{(1)} 
	+ \lambda^2 \dot{a}_m^{(2)} + \cdots = } 	\\ 
&& - \frac{i}{\hbar} \sum_n \left( a_n^{(0)} + \lambda a_n^{(1)} 
  	+ \lambda^2 a_n^{(2)} + \cdots \right) \lambda V_{mn} 
  	e^{-i \omega_{mn} t} \,. \nonumber \\
\end{eqnarray}
%% 
Aqu\'i $\lambda$ es un par\'ametro que luego descartamos, pero nos ayuda 
a indicar que la serie debe converger. Igualando t\'erminos con las mismas 
potencias de $\lambda$ genera el conjunto de ecuaciones
%
\begin{subequations}
\begin{eqnarray} 
\lambda^0 : \qquad \dot{a}_m^{(0)} &=& 0 	\,, 	\\
\lambda^1 : \qquad \dot{a}_m^{(1)} 
 	&=& - \frac{i}{\hbar} \sum_n a_n^{(0)} V_{mn} e^{i \omega_{mn} t} \,, \\ 
\lambda^2 : \qquad \dot{a}_m^{(2)} 
 	&=& - \frac{i}{\hbar} \sum_n a_n^{(1)} V_{mn} e^{i \omega_{mn} t} \,, \\ 
&\vdots& 	\nonumber \\ 
\lambda^r : \qquad \dot{a}_m^{(r)} 
 	&=& - \frac{i}{\hbar} \sum_n a_n^{(r-1)} V_{mn} e^{i \omega_{mn} t} \,. 
\end{eqnarray}
\end{subequations}
%% 

La soluci\'on a orden cero es $a_m^{(0)} = \mathrm{constante}$. Entonces 
los $a_m^{(0)}$ son los valores iniciales del problema, los cuales se 
escogen como $a_{m = i}^{(0)} =1$ y $a_{m \neq i}^{(0)} = 0$, de modo 
que a $t=0$ el sistema ocupa con certeza un estado con energ\'ia $E_i$. 
De esta manera la ecuaci\'on para $a_m^{(1)}$ se reduce a 
%
\begin{eqnarray} 
\dot{a}_m^{(1)} = - \frac{i}{\hbar} V_{mn} e^{i \omega_{mn} t} \,. 
\end{eqnarray}
%% 
$|a_m^{(1)}|^2$ es la (densidad de) probabilidad a primer orden de 
encontrar al sistema a tiempo $t$ en el estado $\ket{m}$, dado que a 
$t=0$ estaba en el estado $\ket{n}$. 

Ahora consideraremos cuatro casos con tratamiento perturbativo a primer 
orden: (i) Excitaci\'on independiente del tiempo $V = V_0$, (ii) excitaci\'on 
arm\'onica $V= V_0 \cos{\nu t}$, (iii) excitaci\'on de un cont\'inuo, 
(iv) excitaci\'on por un campo de banda ancha. 

En ellos suponemos que el sistema se encuentra inicialmente en el estado 
$\Psi (\mbr,0) = u_i (\mbr) = \ket{i}$.  Esto significa que $C_i (0) =1$ y $C_{n\neq i} (0) = 0$. Deseamos conocer las probabilidades de transici\'on 
hacia otros estados. 
 
A primer orden ($n \to i$, $m \to n$), 
%
\begin{eqnarray} 
 \dot{C}_n (t) \simeq   \dot{C}_n^{(1)} (t) 
	= -i \hbar^{-1} V_{ni} e^{i \omega_{ni} t} 	\,.
\end{eqnarray}
%% 


\vspace{3mm}
\textit{Excitaci\'on independiente del tiempo}, $V = V_0$. 
\newline Esta es mas propiamente una excitaci\'on tipo escal\'on debida, por 
ejemplo a la acci\'on de un campo EM est\'atico. Integrando la ecuaci\'on 
anterior,
%
\begin{eqnarray} 
C_n^{(1)} &=& -i \hbar^{-1} \int_0^t V_0 e^{i \omega_{ni} t'} dt' 
	= -i \hbar^{-1} V_{ni} \frac{ e^{i \omega_{ni} t} - 1 }{i \omega_{ni}} 	
	\nonumber \\ 
&=& \frac{1}{i \hbar} V_0  	e^{i \omega_{ni} t/2} 
		\frac{ \sin{\omega_{ni}t/2} }{\omega_{ni}/2} 	\,, 
\end{eqnarray}
%% 
la probabilidad de transici\'on al estado $\ket{n}$ es 
%
\begin{eqnarray} 
 |C_n^{(1)}|^2 = \frac{|V_0|^2 t^2}{\hbar^2} 
	\mathrm{sinc}^2{(\omega_{ni}t/2)} \,,
\end{eqnarray}
%% 
donde $\mathrm{sinc}(x) =\sin{x}/x$ es llamado \textit{seno cociente}.
Tenemos que hacer notar la dependencia cuadr\'atica en el tiempo 
($t^2$). Obviamente, este resultado solo es v\'alido para tiempos cortos,  
donde $\sin{x}/x \sim 1$,  
%
\begin{eqnarray} 
 |C_n^{(1)}|^2 \simeq \frac{|V_0|^2}{\hbar^2} t^2 \,.
 \end{eqnarray}
%% 
Tambi\'en vemos que si las frecuencias de transici\'on $\omega_{ni}$ 
aumentan, la probabilidad de transici\'on disminuye r\'apidamente.
 
Adem\'as, dado que el c\'alculo es perturbativo la probabilidad total de 
transici\'on fuera de $\ket{i}$ debe ser peque\~na,  
%
\begin{eqnarray} 
 P_T = 1 -  |C_i^{(1)}|^2 = \sum_n  |C_n^{(1)}|^2 \ll 1 \,.
\end{eqnarray}
%% 

\vspace{3mm}
\textit{Excitaci\'on arm\'onica}, $V= V_0 \cos{\nu t}$. 
\newline Este es el caso, por ejemplo, de un campo monocrom\'atico de 
frecuencia $\nu$ y amplitud constante,  
%
\begin{eqnarray} 
C_n^{(1)} &=& \frac{V_0}{2i \hbar} \int_0^t  dt' e^{i \omega_{ni} t'} 
	\left( e^{i \nu t'} + e^{- i \nu t'} \right) 		\nonumber \\
&=& \frac{V_0}{2i \hbar } 	\left[ 
	\frac{ e^{i (\omega_{ni} +\nu) t} - 1 }{i (\omega_{ni}  +\nu)} 	
	+ \frac{ e^{i (\omega_{ni} -\nu) t} - 1 }{i (\omega_{ni}  -\nu)} \right] \\ 
\end{eqnarray}
%% 

Si $\omega_{ni} > 0$ y $\omega_{ni} \approx \nu$, el t\'ermino con 
$\omega_{ni} +\nu \approx 2\nu$ tiene oscilaciones r\'apidas que se promedian 
a cero en un tiempo corto, y adem\'as tiene amplitud mucho menor que el 
segundo t\'ermino, por lo que podemos despreciarlo. Esta es la aproximaci\'on 
de onda rotante (RWA -rotating wave approximation), cuya validez se da para 
frecuencias altas (infrarrojo, visible, UV, X). Por esta misma raz\'on, podemos 
ignorar transiciones a niveles con energ\'ias muy diferentes a $\hbar \nu$. Las transiciones ser\'an probables solo si la frecuencia del campo aplicado es muy 
cercana a la frecuencia de la transici\'on. 

Con la RWA se tiene 
%
\begin{eqnarray} 
 |C_n^{(1)}|^2 = \frac{|V_0|^2 }{4\hbar^2} t^2 
	\mathrm{sinc}^2{[(\omega_{ni} -\nu) t/2]} \,.
\end{eqnarray}
%% 

Con excitaci\'on monocrom\'atica es posible limitarse al acoplamiento entre  
dos niveles \'unicamente. Designemos el nivel superior (inferior) con la letra 
$a$ $(b)$, tal que $\ket{\Psi} = C_a(t) e^{-i \omega_a t} \ket{a} 
+ C_b(t) e^{-i \omega_b t} \ket{b}$. Para el \'atomo de hidr\'ogeno la 
transici\'on mas baja es entre los niveles $\ket{b} = \ket{100}$ y 
$\ket{a} = \ket{210}$.  

El Hamiltoniano de la interacci\'on de un dipolo el\'ectrico con un campo 
EM en la aproximaci\'on del dipolo ($r \ll \lambda$) es 
$V = -e \mbr \cdot \mbE (\mbR,t)$, donde $-e \mbr$ es el momento dipolar 
y $\mbR$ es la posici\'on del centro de masa del \'atomo. 
Los elementos diagonales del operador dipolar se anulan, 
$\bra{j} e \mbr \ket{j} = 0$, puesto que $\mbr$ es un operador 
antisim\'etrico y $| u_j |^2$ es sim\'etrico. Los elementos no diagonales 
son $e r_{ab} = ed$. Ignorando (por simplicidad) la dependencia 
espacial  y usando la RWA tenemos 
$V_{ab} (t) = -ed \,E = -ed \,E_0 \cos{\nu t}$. 

Si $a$ ($b$) es el estado final (inicial) $n$ ($i$) se tiene la probabilidad de 
que el \'atomo absorba energ\'ia del campo EM (absorci\'on estimulada), 
%
\begin{eqnarray} 
|C_a^{(1)}|^2 = \frac{|edE_0|^2 }{4\hbar^2} t^2  
	\mathrm{sinc}^2{[(\omega -\nu) t/2]} \,. 	\label{eq:abs_pert1}
 \end{eqnarray}
%% 
Por el contrario, si $a=i$ y $b=n$, se tiene el proceso de emisi\'on 
estimulada, encontrando que $ |C_a^{(1)}|^2 =  |C_b^{(1)}|^2$, es decir, 
las probabilidades de absorci\'on y emisi\'on estimulada son iguales, 
constituyendo un ejemplo de reversibilidad microsc\'opica. Por supuesto, 
las probabilidades pueden decaer, lo cual tratamos en nota aparte. 
De nuevo, el resultado es v\'alido \'unicamente para tiempos cortos. 
Por supuesto, las probabilidades pueden decaer, lo cual tratamos en 
nota aparte. 

En resonancia, $\omega-\nu = \Delta =0$, la probabilidad de transici\'on 
es $P_a (t) = |C_a^{(1)}|^2 = (|edE_0|^2 / 4\hbar^2) t^2$. No podemos 
dejar que $t$ sea grande, porque entonces $P_a$ podr\'ia ser mayor que 
uno! Si $t$ no es muy peque\~na necesitar\'iamos ir al siguiente orden 
de perturbaci\'on, que corregir\'ia el resultado a menor orden. Podr\'iamos 
escribir la condici\'on  de tiempos cortos de $P_a(t) \ll 1$ como 
$t \ll 2\hbar/|edE_0|$.  

\vspace{3mm}
\textit{Excitaci\'on arm\'onica de un cont\'inuo.} %%%%%%
\newline En ciertos casos se tiene una transici\'on desde o hacia un continuo 
de niveles (fotoiionizaci\'on, emisi\'on espont\'anea, bandas electr\'onicas, 
bandas fot\'onicas, etc). En estos casos el cont\'inuo est\'a descrito por una 
distribuci\'on de niveles o densidad de modos $D(\omega)$ [ver, p. ej., ley 
de Planck] y las frecuencias de transici\'on est\'an tan cercanas que las 
cambiamos por una variable cont\'inua, $\omega_{ni} \to \omega$. 
Escribimos la probabilidad de transici\'on como 
%
\begin{eqnarray} 
P_T &\simeq& \sum_{n \neq i}  |C_n^{(1)}|^2 	\rightarrow 
	\int D(\omega) |C^{(1)} (\omega)|^2 d\omega 	\nonumber \\ 
&=& \int d\omega D(\omega) \frac{|V (\omega)|^2}{4\hbar^2} t^2 
	\mathrm{sinc}^2{[(\omega -\nu) t/2]}  		\\ 
	&=&  \int d\omega F(\omega) G(\omega) \,, 	\nonumber 
\end{eqnarray}
%% 
donde 
%
\begin{eqnarray} 
F(\omega) &=& D(\omega) \frac{|V (\omega)|^2}{4\hbar^2} \,, 
	\qquad G(\omega) = t^2 \mathrm{sinc}^2{[(\omega -\nu) t/2]} \,. 
	\nonumber \\
\end{eqnarray}
%% 
  
Ocurre con frecuencia que $D(\omega)$ y $|C(\omega)|^2$ tienen escalas 
frecuenciales muy distintas, en las que la funci\'on $F(\omega)$ var\'ia muy 
poco comparado con $G(\omega)$,  entonces 
%
\begin{eqnarray} 
J = \int d\omega F(\omega) G(\omega) \sim 
	F(\omega_0) \int G(\omega) d\omega \,.
\end{eqnarray}
%% 
((Un comportamiento an\'alogo ocurre en el dominio temporal, lo que 
lleva al m\'etodo de eliminaci\'on adiab\'atica, que es muy \'util para 
resolver ecuaciones acopladas.)

Por lo general es muy dif\'icil resolver la integral de $P_T$ en forma 
anal\'itica, por lo que requerimos aproximaciones. Para tiempos muy 
cortos tales que $|\omega_{ni}|  t \ll 1 \to \mathrm{sinc}{(x)} \to 1$, 
%
\begin{eqnarray} 
P_T &\simeq& \frac{t^2}{4\hbar^2} \int D(\omega) |V (\omega)|^2 d\omega 	
\end{eqnarray}
%% 
de modo que la raz\'on de transici\'on $\dot{P}_T \propto t$, que implica 
que el sistema tarda en responder (p. ej., un fotodetector) a menos que 
$D(\omega) |V (\omega)|^2$ sea muy ancha. 

Para tiempos mas largos $F(\omega)$ var\'ia muy poco en el rango de 
$G(\omega)$. Entonces $F(\omega) \to F(\nu)$ y 
%
\begin{eqnarray} 
P_T &=& \frac{D(\nu) |V(\nu)|^2}{4\hbar^2} t^2 \int 
	  \mathrm{sinc}^2{[(\omega -\nu) t/2]}  d\omega		\\ 
	x &=& (\omega -\nu)/2 	\qquad \to dx = t d\omega/2 	\\ 
P_T &=& \frac{D(\nu) |V(\nu)|^2}{2\hbar^2} t 
	\int_{-\infty}^{\infty}    \mathrm{sinc}^2{x}  \,dx 	\nonumber \\ 
P_T &=& \frac{\pi}{2\hbar^2} D(\nu) |V(\nu)|^2 t 		\,.
\end{eqnarray}
%% 
Extendimos los l\'imites de la integral a $\pm \infty$ puesto que 
$\mathrm{sinc}^2{x} \ll 1$ para $x \gg 1$. Es decir, las diferencias de 
frecuencias que son grandes  casi no aportan a la integral.La raz\'on de 
transici\'on es, entonces 
% 
\begin{eqnarray} 
\Gamma &=& \dot{P}_T = - \partial_t |C_i^{(1)}|^2 	\nonumber \\ 
	&=& \frac{\pi}{2\hbar^2} D(\nu) |V(\nu)|^2 = \mathrm{const.}		\,,
\end{eqnarray}
%% 
que es la llamada \textit{Regla de Oro} de Fermi. Es proporcional a la 
densidad de modos y a la intensidad del campo aplicado. 

Para tiempos muy largos $P_T$ no es peque\~na y no podemos 
suponer que $|C_i^{(1)}|^2 \simeq 1$. Entonces debemos multiplicar 
$\Gamma$ por $|C_i^{(1)}|^2 \sim 1$, con lo que generalizamos la 
Regla de Oro de Fermi: $\Gamma |C_i|^2 = - \partial_t |C_i|^2$,  y 
obtenemos el comportamiento t\'ipico de decaimiento exponencial 
$|C_i|^2 \propto e^{- \Gamma t}$. 

\vspace{3mm}
\textit{Excitaci\'on por luz de banda ancha.}
\newline Un caso similar es el de excitaci\'on de un \'atomo de dos niveles 
por un campo no monocrom\'atico, como el de radiaci\'on de un cuerpo negro, 
cuyas frecuencias $\nu$ forman un cont\'inuo. En este caso $V = V(\nu)$, y 
la probabilidad total de transici\'on toma la forma  
%
\begin{eqnarray}  
P_T &=& \int d\nu \mathcal{U}(\nu) \frac{t^2}{\hbar^2}   
	  \mathrm{sinc}^2{[(\omega -\nu) t/2]}   		\,,
\end{eqnarray}
%% 
donde $\mathcal{U}(\nu) \propto E_0^2$ es la densidad espectral de 
energ\'ia y la integral es sobre todas las frecuencias $\nu$. Ahora vemos que 
$\mathcal{U}(\nu)$ es una curva mucho m\'as ancha que $(\sin{x}/x)^2$, por 
lo que la sacamos de la integral como $\mathcal{U}(\omega)$. La raz\'on de 
transici\'on $\dot{P}_T$ es 
%
\begin{eqnarray}  
B(\omega) \mathcal{U}(\omega) 
	= \frac{\pi}{3\hbar^2 \epsilon_0} \wp^2 \mathcal{U}(\omega) \,,
\end{eqnarray}
%% 
donde el 3 surge de reemplazar $\wp^2$ por $\wp^2/3$ (  $\wp$ es el 
momento dipolar del \'atomo) por que la radiaci\'on viene de todas las 
direcciones, y $B(\omega)$ es el coeficiente de absorci\'on y emisi\'on 
estimulada. 

Ver apuntes de F\'isica de L\'aseres. 

%%%%%%%%%%%%%%%%%%%%%%%%%%%%%%%%%%%
%%%%%%%%%%%%%%%%%%%%%%%%%%%%%%%%%%%
\subsection{Excitaci\'on Monocrom\'atica: F\'ormula de Rabi} 
El ancho de banda de algunos l\'aseres es muy angosto, lo que permite 
usarlos para excitar selectivamente transiciones at\'omicas, es decir, sin 
excitar a su vez otras transiciones. En un curso de l\'aseres aprendemos 
que el n\'umero de fotones en un ancho de banda angosto puede ser muy 
grande, de tal manera que el tratamiento perturbativo se vuelve engorroso si 
tuvi\'eramos que calcular la probabilidad de transici\'on a varios \'ordenes 
debido a que la excitaci\'on no sea d\'ebil. Afortunadamente, hay maneras 
m\'as sencillas de realizar c\'alculos exactos, o muy aproximados, que son 
v\'alidos para $\Omega$ arbitraria (pero $\Omega$ mucho menor que la 
frecuencia de transici\'on, algo que se cumple en much\'isimos casos). En 
esta secci\'on y las siguientes trataremos el tema, sus resultados principales 
y m\'etodos alternativos, en la situaci\'on en la que se ignoran las p\'erdidas 
(de energ\'ia y polarizaci\'on). Posteriormente tomaremos en cuenta estos 
efectos.  

Consideraremos un campo EM monocrom\'atico de frecuencia $\nu$ cerca de resonancia con la transici\'on de frecuencia $\omega =\omega_a - \omega_b$. 
Los niveles superior e inferior de la transici\'on est\'an descritos por los estados 
$\ket{a}$ y $\ket{b}$, respectivamente. 

Adelantamos la aplicaci\'on de la RWA escribiendo el Hamiltoniano de la 
interacci\'on como ($V_{mn} = \bra{m} V \ket{n}$)
%
\begin{eqnarray*} 
V_{ab} &=& -d E_0 \cos{\nu t} 	\simeq 	- (1/2) d E_0 e^{-i \nu t} 	\\
V_{ba} &=& -d E_0 \cos{\nu t} 	\simeq 	- (1/2) d E_0 e^{i \nu t} \,. 
\end{eqnarray*}
%% 
Las ecuaciones para las amplitudes se obtienen de (\ref{eq:amplitudeEqsIP}) 
y son 
%
\begin{subequations}
\begin{eqnarray} 
\dot{C}_a (t) &=& - \frac{i}{\hbar} \bra{a} V_0 \ket{b} 
    e^{-i \nu t} e^{i \omega t} C_b (t)
	= i \frac{\mu_0}{2} e^{i \Delta t} C_b (t) 	\label{eq:amp_a1} \nonumber \\ \\ 
\dot{C}_b (t) &=& - \frac{i}{\hbar} \bra{b} V_0 \ket{a} e^{i \nu t} 
	e^{-i \omega t} C_a (t)
	= i \frac{\mu_0}{2} e^{-i \Delta t} C_a (t) 	\label{eq:amp_b1} \nonumber \\ 
\end{eqnarray}
\end{subequations}
%% 
donde $\mu_0 = d E_0 /\hbar$ y  $\Delta = \omega -\nu$. 

Para resolver estas ecuaciones acopladas supongamos una soluci\'on de la 
forma $C_b = e^{i \mu t} \to \dot{C}_b = i \mu e^{i \mu t}$, que insertamos en 
la Ec. (\ref{eq:amp_b1}):  
%
\begin{eqnarray} 
C_a (t) &=& \frac{2 \mu}{\mu_0} e^{i (\Delta +\mu)t} \,, \label{eq:amp_a2}  
\end{eqnarray}
%% 
que insertamos a su vez en la Ec. (\ref{eq:amp_a1})
%
\begin{eqnarray*} 
i \frac{2 \mu}{\mu_0} (\mu +\Delta) e^{i (\mu + \Delta)t} 
	&=& i \frac{\mu_0}{2} e^{i (\mu + \Delta)t} 	\,, 	\\ 
\mu^2 + \Delta \mu - \frac{\mu_o^2}{4} &=& 0 	\,.
\end{eqnarray*}
%% 
Las soluciones de esta ecuaci\'on cuadr\'atica son  
%
\begin{eqnarray} 
\mu_{\pm} &=& -\frac{\Delta}{2} \pm \frac{1}{2} \sqrt{\mu_0^2 +\Delta^2} \,. 
\label{eq:mu_pm}  
\end{eqnarray}
%% 
Las soluciones para $C_b$ y $C_a$ [habiendo sustitu\'ido $C_b$ en la 
Ec.(\ref{eq:amp_b1})] toman las formas 
%
\begin{subequations}
\begin{eqnarray} 
C_b (t) &=& A e^{i \mu_+ t} + B e^{i\mu_- t}  \,, \\ %\label{eq:amp_b3}  
C_a (t) &=& \frac{2}{\mu_0} e^{i \Delta t} 
	\left[ A \mu_+ e^{i \mu_+ t} + B \mu_- e^{i \mu_- t} \right] \,. 
\end{eqnarray}
\end{subequations}
%% 

Para determinar las constantes $A$ y $B$ requerimos condiciones iniciales. 
Para ello consideremos que el \'atomo inicia en su estado base, $C_b(0) =1$ 
y $C_a (0) =0$, lo que significa que calcularemos la probabilidad de absorci\'on:  
%
\begin{eqnarray*} 
C_b (0) &=& 1 = A + B   \,, 	\\  
C_a (0) &=& 0 = \frac{2}{\mu_0} \left[ A \mu_+  + B \mu_-  \right] \,. 	 
\end{eqnarray*}
%% 
Redefiniendo $\mu_0 = \Omega_0$ y usando 
$\mu_+ \mu_- = - \frac{1}{4} \Omega_0^2$ se tiene 
%
\begin{eqnarray} 
A &=& - \frac{\mu_-}{\Omega} \,, 	\qquad B = \frac{\mu_+}{\Omega} \,, 
	\nonumber \\ 
\Omega &=& \mu_+ - \mu_-    = \sqrt{\Omega_0 +\Delta^2}  \,. 
\end{eqnarray}
%% 
Finalmente, la probabilidad de absorci\'on se obtiene como 
%
\begin{eqnarray} 
C_a (t) &=& \frac{2}{\Omega_0} e^{i \Delta t} 
	\left[ - \frac{\mu_- \mu_+}{\Omega}  e^{-i \Delta t/2}  e^{i \Omega t/2} 
	\right. 	\nonumber \\ 
	&& \left. + \frac{\mu_- \mu_+}{\Omega}  e^{-i \Delta t/2}  e^{-i \Omega t/2}
	 \right]  	\nonumber \\ 
&=& \frac{\Omega_0}{2 \Omega} e^{i \Delta t/2} 
	\left[ e^{i \Omega t/2} - e^{-i \Omega t/2} \right] 	\nonumber \\ 
&=& i \frac{\Omega_0}{\Omega} e^{i \Delta t/2}  \sin{(\Omega t/2)} \,, 
	\nonumber \\
|C_a (t)|^2 &=& \frac{\Omega_0^2}{\Omega^2} \sin^2{(\Omega t/2)}  \,. 
\end{eqnarray}
%% 
Esta expresi\'on es llamada la \textit{F\'ormula de Rabi} y caracteriza la 
interacci\'on coherente entre el \l\'aser y el \'atomo, en la que el primero 
imparte su coherencia al segundo). Las oscilaciones propias de la 
probabilidad de transici\'on ocurren a la frecuencia $\Omega/2$, donde 
$\Omega_0 \equiv \wp E_0 /\hbar$ es la llamada \textit{Frecuencia de Rabi} 
y $\Omega$ es una frecuencia de Rabi generalizada. Al estudiar procesos 
de interacci\'on entre dos sistemas, las oscilaciones indican que hay 
coherencia. 

Comparemos la F\'ormula de Rabi con el caso perturbativo, 
Ec.(\ref{eq:abs_pert1}), en el que la frecuencia de las oscilaciones ocurre esta 
dada por la desinton\'ia entre las frecuencias del l\'aser y de la transici\'on. 
Adem\'as, el factor $\Omega^{-2}$, dependiente de $\Omega_0$, indica dos 
cosas: (i) la probabilidad se ensancha por potencia y (ii) la soluci\'on tambi\'en 
est\'a permitida en resonancia, $\Delta =0$ (el denominador no es cero en 
este caso), en cuyo caso $|C_a (t)|^2 = \sin^2{(\Omega t/2)}$ (oscilaciones 
entre cero y uno). 

La probabilidad de permanecer en el estado base se obtiene como: 
%
\begin{eqnarray} 
C_b (t) &=& e^{i \Delta t/2} 
	\left[ - \frac{\mu_-}{\Omega}  e^{i \Omega t/2} 
	 + \frac{\mu_+}{\Omega}   e^{-i \Omega t/2} \right]  	\nonumber \\ 
&=& e^{i \Delta t/2} \left[ \frac{\Delta +\Omega}{2\Omega} e^{i \Omega t/2} 
	 - \frac{\Delta -\Omega}{2\Omega}   e^{-i \Omega t/2} \right] 
	 \nonumber \\ 
&=& e^{i \Delta t/2} \left[ \cos{(\Omega t/2)} 
	+ i \frac{\Delta}{\Omega} \sin{(\Omega t/2)} \right] \,, \nonumber \\ 
|C_b (t)|^2 &=& \cos^2{(\Omega t/2)} 
	+ \frac{\Delta^2}{\Omega^2} \sin^2{(\Omega t/2)}  \,.
\end{eqnarray}
%% 

%%5%%%%%%%%60%%%%%%%%70%%%
\subsection{Descripciones Matricial y de Operadores de Proyecci\'on}
Consideremos el problema anterior desde una perspectiva mas algebraica, 
que permitir\'a un tratamiento mas completo y a\~nadir mas resultados. 

Primero definimos las matrices (operadores) de esp\'in de Pauli y la matriz 
identidad $2 \times 2$:
%
\begin{eqnarray*}
\sigma_x &=& \left( \begin{array}{cc}  0&1 \\ 1&0 \end{array} \right) \,, 
\quad \sigma_y = \left( \begin{array}{cc} 
	0&-i \\ i&0 \end{array} \right) \,, 	\quad 
\sigma_z =\left( \begin{array}{cc} 1&0 \\ 0&-1 \end{array} \right) \,, \\
\sigma_+ &=& \left( \begin{array}{cc} 0&1 \\ 0&0 \end{array} \right) \,, 
\quad \sigma_- = \left( \begin{array}{cc} 
	0&0 \\ 1&0 \end{array} \right) \,, 
\qquad \mathbf{1} = \left( \begin{array}{cc} 
	1&0 \\ 0&1 \end{array} \right) \,,
\end{eqnarray*}
%%
y los vectores 
%
\begin{eqnarray*}
\ket{a} &=& \left( \begin{array}{c}  1 \\ 0 \end{array} \right) \,, \qquad 
\ket{b} = \left( \begin{array}{c}  0 \\ 1 \end{array} \right) \,, 
\end{eqnarray*}
%%
Para una notaci\'on m\'as compacta podemos usar operadores de proyecci\'on 
en lugar de matrices: 
%
\begin{eqnarray*}
\sigma_+ &=& \ketbra{a}{b} \,, 	\qquad 
	\sigma_- = \ketbra{b}{a} \,, 	\\ 
\sigma_z &=& [\sigma_+, \sigma_-]   
	= \ketbra{a}{a} - \ketbra{b}{b} 	\,, \\ 
\mathbf{1} &=& \ketbra{a}{a} + \ketbra{b}{b} \,.	
\end{eqnarray*}
%%
Esta \'ultima es la relaci\'on de completez.  

Ahora escribimos la funci\'on de onda del sistema de dos niveles como 
$\ket{\Psi (t)} = C_a (t) \ket{a} + C_b (t) \ket{b}$, y establecemos 
el origen de la energ\'ia del sistema de dos niveles justo a la mitad, 
tal que $\omega_a = - \omega_b = \omega /2$. Utilizando 
$H_0 = \mathbf{1} H_0 \mathbf{1}$, $H_0 \ket{k} = \hbar \omega_k \ket{k}$ 
y $\braket{m}{n} = \delta_{mn}$, el Hamiltoniano de la energ\'ia del \'atomo 
se puede escribir como
%
\begin{eqnarray*} 
H_0 &=& (\ketbra{a}{a} + \ketbra{b}{b}) H_0 
 	( \ketbra{a}{a} + \ketbra{b}{b} ) 		\\ 
&=& (\ketbra{a}{a} + \ketbra{b}{b}) 
 ( \hbar \omega_a  \ketbra{a}{a} + \hbar \omega_b  \ketbra{b}{b} ) \\ 
&=& \hbar \omega_a  \ketbra{a}{a} 
    +\hbar \omega_b  \ketbra{b}{b} \\  
&=& (\hbar \omega /2)( \ketbra{a}{a} - \ketbra{b}{b} ) 
	= (\hbar \omega /2) \sigma_z
\end{eqnarray*}
%% 

El operador de la interacci\'on por acoplamiento dipolar es 
$V(t) = -e \mbr \cdot \mbE (R,t)$. Por simplicidad 
consideramos que el campo est\'a linealmente polarizado en la 
direcci\'on $x$, entonces 
$V(t) = -e x E_x \cos{\nu t} = - \wp E_x \cos{\nu t} $, donde 
%
\begin{eqnarray*}
\wp &=& \wp_{ab} \ketbra{a}{b} +\wp_{ba} \ketbra{b}{a} 
  = \left( \begin{array}{cc} 0& \wp_{ab}  \\ \wp_{ba} &0 
\end{array} \right) \,.  
\end{eqnarray*}
%%
Similarmente, $V(t)$ se puede escribir como
%
\begin{eqnarray*} 
V(t) &=& - E_x(t) (\ketbra{a}{a} + \ketbra{b}{b}) \wp 
 	( \ketbra{a}{a} + \ketbra{b}{b} ) 		\\ 
&=& -E_x(t) (\wp_{ab} \ketbra{a}{b} + \wp_{ba} \ketbra{b}{a} ) \\  
&\simeq& - \frac{E_0}{2}  \left( \wp_{ab} e^{-i \nu t} \ketbra{a}{b}  
	+ \wp_{ba} e^{i \nu t} \ketbra{b}{a}  \right) \,,
\end{eqnarray*}
%% 
donde hemos aplicado la RWA (ignoramos los t\'erminos que oscilan a 
frecuencias $\omega + \nu$) considerando la evoluci\'on libre del dipolo, 
donde $\wp_{ab}$ contiene impl\'icitamente el factor $e^{i \omega t}$ y 
$\wp_{ba}$ el factor $e^{i \omega t}$.  El Hamiltoniano completo 
toma la forma  
%
\begin{eqnarray*}
H &=& \frac{1}{2} 
\left( \begin{array}{cc}  \hbar \omega & -\wp_{ab} E_0 e^{-i \nu t} \\ 
	- \wp_{ba} E_0 e^{i \nu t} &- \hbar \omega  \end{array} \right) 	\\ 
&=& \frac{\hbar}{2} 
\left( \begin{array}{cc}  \omega & - \Omega_0 e^{-i \nu t} \\ 
	- \Omega_0 e^{i \nu t} & - \omega  \end{array} \right) \,.
\end{eqnarray*}
%%

Con esto la ecuaci\'on de Schr\"odinger 
$\ket{\dot{\Psi}} = -(i/\hbar) \ket{\Psi}$ se escribe como 
%
\begin{eqnarray*}
\left( \begin{array}{c} \dot{C}_a \\ \dot{C_b} \end{array}  \right)
	&=& \frac{i}{2} 
\left( \begin{array}{cc}  - \omega & \Omega_0 e^{-i \nu t} \\ 
	 \Omega_0 e^{i \nu t} & \omega  \end{array} \right) 
	 \left( \begin{array}{c} C_a \\ C_b \end{array}  \right) \,.
\end{eqnarray*}
%%
Haciendo los cambios de variables $C_a(t) = c_a(t) e^{-i \nu t/2}$ y 
$C_b(t) = c_b(t) e^{i \nu t/2}$ (cambio de la imagen de Schr\"odinger a la de 
Interacci\'on) tenemos finalmente 
%
\begin{eqnarray}
\left( \begin{array}{c} \dot{c}_a \\ \dot{c}_b \end{array}  \right)
	&=& \frac{i}{2} 
\left( \begin{array}{cc}  - \Delta & \Omega_0  \\ 
	 \Omega_0  & \Delta  \end{array} \right) 
	 \left( \begin{array}{c} c_a \\ c_b \end{array}  \right) \,.
\end{eqnarray}
%%
Esto es completamente equivalente al m\'etodo visto anteriormente. 

%%%%%%%%%%%%%%
\subsection{Estados Vestidos Semicl\'asicos \label{dressedST}}  
Una descripci\'on alternativa de la interacci\'on \'atomo-l\'aser de gran 
importancia es la de usar los estados vestidos, que son los vectores propios 
de la matriz en las ecuaciones para las amplitudes en lugar de los 
estados del \'atomo. Obtendremos dichos vectores propios y la matriz de transformaci\'on entre las bases. 

Tenemos que resolver la ecuaci\'on 
%
\begin{eqnarray*}
\mathbf{M} \left[ \begin{array}{c} u \\ v \end{array} \right] 
    = \left[ \begin{array}{cc} -\delta & \Omega_0 \\ 
    \Omega_0 & \delta \end{array} \right]  
    \left[ \begin{array}{c} u \\ v \end{array} \right] 
    = \lambda \left[ \begin{array}{c} u \\ v \end{array} \right]
\end{eqnarray*}
%% 
Valores propios ($\det(\mathbf{M} -\lambda) =0$):  
%
\begin{eqnarray*}
\left| \begin{array}{cc} -\delta - \lambda & \Omega_0 \\ 
    \Omega_0 & \delta -\lambda \end{array} \right| 
    &=& -(\delta +\lambda)(\delta -\lambda) - \Omega_0^2 = 0  \\  
\to \lambda^2 -\delta^2 - \Omega_0^2 &=& 0 \quad 
    \to \lambda_{\pm} \pm \Omega = \pm \sqrt{\Omega_0^2 +\delta^2} .
\end{eqnarray*}
%% 
La separaci\'on de energ\'ia entre los estados es 
$\Delta E = 2\hbar \Omega$. 

El vector propio para $\lambda_-$: 
%
\begin{eqnarray*}
\left[ \begin{array}{cc} -\delta +\Omega & \Omega_0 \\ 
    \Omega_0 & \delta +\Omega \end{array} \right]  
    \left[ \begin{array}{c} u_1 \\ v_1 \end{array} \right] &=& 0 \\ 
%
(\Omega -\delta) u_1 + \Omega_0 v_1 &=& 0     \\ 
\Omega_0 u_1 +(\delta +\Omega) v_1 &=& 0 
\end{eqnarray*}
%% 
Hacemos $v_1 =1 \to u_1 = -\Omega_0/(\Omega -\delta)$. 

Vector propio para $\lambda_+$: 
%
\begin{eqnarray*}
\left[ \begin{array}{cc} -\delta -\Omega & \Omega_0 \\ 
    \Omega_0 & \delta -\Omega \end{array} \right]  
    \left[ \begin{array}{c} u_2 \\ v_2 \end{array} \right] &=& 0 \\ 
%
-(\Omega +\delta) u_2 + \Omega_0 v_2 &=& 0     \\ 
\Omega_0 u_2 +(\delta -\Omega) v_2 &=& 0 
\end{eqnarray*}
%% 
Hacemos $u_1 =1 \to v_1 = \Omega_0/(\Omega -\delta)$. 

El factor de normalizaci\'on es $$N_j^2 = u_j^2 +v_j^2 
= 1 + \frac{\Omega_0^2}{(\Omega -\delta)^2} 
  = \frac{(\Omega -\delta)^2 +\Omega_0^2}{(\Omega -\delta)^2}$$
por lo que los vectores normalizados son: 
%
\begin{eqnarray*}
\ket{-} &=& \left[ \begin{array}{c} u_1 \\ v_1 \end{array} \right] 
    = \frac{1}{ \sqrt{(\Omega -\delta)^2 +\Omega_0^2} } 
    \left[ \begin{array}{c} -\Omega_0 \\ \Omega -\delta 
    \end{array} \right] 
    = \left[ \begin{array}{c} -\sin{\theta} \\ \cos{\theta}  
    \end{array} \right]  \\ 
    &=& -\sin{\theta} \ket{e} + \cos{\theta} \ket{g}    \\
%
\ket{+} &=& \left[ \begin{array}{c} u_1 \\ v_1 \end{array} \right] 
    = \frac{1}{ \sqrt{(\Omega -\delta)^2 +\Omega_0^2} } 
    \left[ \begin{array}{c} \Omega -\delta \\ \Omega_0  
    \end{array} \right] 
    = \left[ \begin{array}{c} \cos{\theta} \\ \sin{\theta}  
    \end{array} \right] \\ 
    &=& \cos{\theta} \ket{e} + \sin{\theta} \ket{g}    \\
\end{eqnarray*}
%% 

La matriz de  transformaci\'on es 
%
\begin{eqnarray*}
U &=& \left[ \begin{array}{cc} u_2 & v_2 \\ 
    u_1 & v_1 \end{array} \right]  
    = \left[ \begin{array}{cc} \cos{\theta} & \sin{\theta} \\ 
    -\sin{\theta} & \cos{\theta} \end{array} \right]     \\ 
U^{-1} &=& \left[ \begin{array}{cc} u_2 & u_1 \\ 
    v_2 & v_1 \end{array} \right]  
    = \left[ \begin{array}{cc} \cos{\theta} & -\sin{\theta} \\ 
    \sin{\theta} & \cos{\theta} \end{array} \right] . 
\end{eqnarray*}
%% 

%
\begin{eqnarray*}
\sin{2\theta} &=& 2 \cos{\theta} \,\sin{\theta} 
 = \frac{2(\Omega -\delta) \Omega_0}{(\Omega -\delta)^2 +\Omega_0^2} \\ 
&=& \frac{2(\Omega -\delta) \Omega_0}
    {\Omega^2 +\delta^2 -2 \Omega \delta +\Omega_0^2} 
    = \frac{2(\Omega -\delta) \Omega_0}
    {2( \Omega_0^2 +\delta^2 -\Omega \delta )}     \\ 
    &=& \frac{(\Omega -\delta) \Omega_0}{\Omega^2 -\delta \Omega} 
    = \frac{\Omega_0}{\Omega} .
\end{eqnarray*}
%%
%
\begin{eqnarray*}
\cos{2\theta} &=& \cos^2{\theta} -\sin^2{\theta} 
 = \frac{(\Omega -\delta)^2 -\Omega_0^2}{(\Omega -\delta)^2 
     +\Omega_0^2} \\ 
&=& \frac{\Omega^2 +\delta^2 -2 \Omega \delta -\Omega_0^2}
    {\Omega^2 +\delta^2 -2 \Omega \delta +\Omega_0^2} 
    = \frac{2(\delta -\Omega) \delta} 
    {2( \Omega^2 -\Omega \delta )}     \\ 
    &=& \frac{-\delta}{\Omega} .
\end{eqnarray*}
%%

Veamos casos especiales de los estados vestidos:  
\newline (i) $\delta =0$, $\cos{\theta} = \sin{\theta} 
=1/\sqrt{2}$, $\Delta E = 2\hbar \Omega_0$. Entonces 
$\ket{\pm} = (\pm \ket{e} +\ket{g} )/ \sqrt{2}$. 
\newline (i) $|\delta| \gg \Omega_0$ y $\delta >0$, 
$\cos{\theta} \simeq 0$, $\sin{\theta} \simeq 1$, 
$\Delta E = 2\hbar |\delta|$. Entonces 
$\ket{+} \simeq \ket{g}$ y $\ket{-} \simeq -\ket{e}$. 
\newline (i) $|\delta| \ll \Omega_0$ y $\delta <0$, 
$\cos{\theta} \simeq 1$, $\sin{\theta} \simeq 0$, 
$\Delta E = 2\hbar |\delta|$. Entonces 
$\ket{+} \simeq \ket{e}$ y $\ket{-} \simeq \ket{g}$. 
 
Los estados vestidos definen la matriz de transformaci\'on 
$U$ que diagonaliza la matriz $\mathbf{M}$ como 
%
\begin{eqnarray*}
U \mathbf{M} U^{-1} = \left[ \begin{array}{cc} 
    \Omega & 0 \\ 0 & -\Omega \end{array} \right] \,.
\end{eqnarray*}
%%
Esta matriz de transformaci\'on relaciona las amplitudes de 
probabilidad de los estados vestidos con las de los estados 
de la base original (base desnuda) por 
%
\begin{eqnarray*}
\left[ \begin{array}{c} C_2(t) \\ C_1(t) \end{array} \right] 
= U \left[ \begin{array}{c} C_a(t) \\ C_b(t) \end{array} \right]    \,.
\end{eqnarray*}
%%
donde las amplitudes de probabilidad $C_1,\,C_2$ obedecen las 
ecuaciones de movimiento  (proof?) 
%
\begin{eqnarray*}
\dot{C}_1 &=& (-i/2) \Omega C_1 \,, 
    \qquad \to C_1(t) = C_1^0 e^{-i \Omega t/2} \,,\\ 
\dot{C}_2 &=& (i/2) \Omega C_2 \,, 
    \qquad \quad \to C_2(t) = C_2^0 e^{i \Omega t/2} \,.  
\end{eqnarray*}
%%

Obtendremos las soluciones de $[ C_a(t), \, C_b(t) ]$ en 
t\'erminos de sus condiciones iniciales. Escribimos los valores 
iniciales $[ C_2^0, \, C_1^0 ]$ en t\'erminos de 
$[ C_a^0, \, C_b^0 ]$, 
%
\begin{eqnarray*}
\left[ \begin{array}{c} C_2^0 \\ C_1^0 \end{array} \right] 
&=& U \left[ \begin{array}{c} C_a^0 \\ C_b^0 \end{array} \right] 
= \left[ \begin{array}{c} C_a^0\cos{\theta} + C_b^0 \sin{\theta} \\ 
    -C_a^0\sin{\theta} + C_b^0\cos{\theta} \end{array} \right] 
\end{eqnarray*}
%%
%
\begin{eqnarray*}
\left[ \begin{array}{c} C_2(t) \\ C_1(t) \end{array} \right] 
&=& \left[ \begin{array}{c} e^{i\Omega t/2}
    ( C_a^0\cos{\theta} + C_b^0 \sin{\theta} ) \\ 
    e^{-i\Omega t/2}
    (-C_a^0\sin{\theta} + C_b^0\cos{\theta} ) \end{array} \right] 
\end{eqnarray*}
%%
Entonces 
\begin{widetext}
%
\begin{eqnarray*}
\left[ \begin{array}{c} C_a(t) \\ C_b(t) \end{array} \right]    
&=& U^{-1} 
\left[ \begin{array}{c} C_2(t) \\ C_1(t) \end{array} \right]  
%
  = \left[ \begin{array}{c} e^{i\Omega t/2}
    ( C_a^0\cos^2{\theta} + C_b^0 \sin{\theta} \cos{\theta} ) 
    + e^{-i\Omega t/2}
    (C_a^0\sin^2{\theta} - C_b^0 \sin{\theta} \cos{\theta} ) \\ 
e^{i\Omega t/2}
    ( C_a^0 \sin{\theta} \cos{\theta} + C_b^0 \sin^2{\theta} ) 
    + e^{-i\Omega t/2}
    (-C_a^0 \sin{\theta} \cos{\theta} + C_b^0 \cos^2{\theta} )
    \end{array} \right]     \\ 
%
&=& \left[ \begin{array}{c} e^{i\Omega t/2}
 [ ((\Omega -\delta)/2\Omega) C_a^0 + (\Omega_0/2\Omega) C_b^0  ] 
 + e^{-i\Omega t/2} 
 [ ((\Omega +\delta)/2\Omega) C_a^0 - (\Omega_0/2\Omega) C_b^0  ] \\ 
e^{i\Omega t/2} 
    [ (\Omega_0/2\Omega) C_a^0 + ((\Omega +\delta)/2\Omega) C_b^0 ] 
    + e^{-i\Omega t/2}
    [ (-\Omega_0/2\Omega) C_a^0 + ((\Omega -\delta)/2\Omega) C_b^0 ]
    \end{array} \right]      \\
%
&=& \left[ \begin{array}{c} 
C_a^0 [ \cos{(\Omega t/2)} -(i\delta /\Omega) \sin{(\Omega t/2)} ] 
    + (i\Omega_0 /\Omega) C_b^0 \sin{(\Omega t/2)}     \\ 
    (i\Omega_0 /\Omega) C_a^0 \sin{(\Omega t/2)} 
    +C_b^0 
    [ \cos{(\Omega t/2)} +(i\delta /\Omega) \sin{(\Omega t/2)} ]
 \end{array} \right]         \\ 
%
&=& \left[ \begin{array}{c} 
 [ \cos{(\Omega t/2)} -(i\delta /\Omega) \sin{(\Omega t/2)} ] 
    + (i\Omega_0 /\Omega)  \sin{(\Omega t/2)}     \\ 
    (i\Omega_0 /\Omega) \sin{(\Omega t/2)} 
    +[ \cos{(\Omega t/2)} +(i\delta /\Omega) \sin{(\Omega t/2)} ]
 \end{array} \right]  
 \left[ \begin{array}{c} C_a^0 \\ C_b^0 \end{array} \right] \,.
\end{eqnarray*}
%%

\end{widetext} 


%%%%%%%%%%%%%%%%%%%%%%%%%%%%%%%%%%%
%%%%%%%%%%%%%%%%%%%%%%%%%%%%%%%%%%%
\section{Operador de Densidad} 
La adici\'on fenomenol\'ogica de disipaci\'on no es correcta en muchas 
situaciones, por lo que se recurre al operador de densidad. T\'ipicamente, 
no conocemos el estado completo del vector para un sistema de muchas 
part\'iculas. Cada \'atomo tiene una historia \'unica de emisiones, absorciones, colisiones, etc., pero conocemos ciertas propiedades estad\'isticas. El 
operador de densidad contiene esa informaci\'on. Definimos el operador de 
densidad como 
%
\begin{eqnarray}
\rho &=& \overline{ \ketbra{\Psi}{\Psi} }	
	= \sum_{\Psi} P_{\Psi} \ketbra{\Psi}{\Psi} 	\,.
\end{eqnarray}
%%
La suma sobre $\Psi$ representa promedios cl\'asicos sobre el ensamble, 
y puede implicar varias sumas e integrales. $P_{\Psi}$ es la fracci\'on de 
sistemas con la funci\'on de onda $\ket{\Psi}$.  Para vectores de la forma 
$\ket{\Psi} = \sum_n C_n \ket{n}$, el operador de densidad es 
%
\begin{eqnarray*}
\rho &=& \sum_{\Psi} P_{\Psi} \sum_n \sum_m C_n C_m^\ast 
	\ketbra{n}{m} 	\\ 
&=& \sum_n \sum_m \rho_{nm} \ketbra{n}{m}  	\,.
\end{eqnarray*}
%%

\vspace{2mm} %%%%%%%%%%%%%%%%%
%\begin{widetext}
\textbf{Ejemplo: Meystre +Sargent 4.2} Considere una colecci\'on de 
\'atomos de dos niveles de los cuales 30\% est\'an descritos por la
 funci\'on de onda $\ket{\Psi_1(t)} = 2^{-1/2} [e^{-i \omega_a t} \ket{a} 
+e^{-i \omega_b t} \ket{b} ]$, 50\% descritos por 
$\ket{\Psi_2(t)} = 10^{-1/2} [ e^{-i \omega_a t} \ket{a} 
-3e^{-i \omega_b t} \ket{b} ]$, y 20\% descritos por 
$\ket{\Psi_3(t)} = e^{-i \omega_b t} \ket{b}$. ?`Cu\'al es la probabilidad 
de que este sistema est\'e en el estado $\ket{\Psi_j}$? Demuestre que 
$\rho^2 \neq \rho$. 

\textit{Soluci\'on}
%
\begin{eqnarray*}
\ket{\Psi_1} &=& \frac{1}{\sqrt{2}} \left( \begin{array}{c} e^{-i \omega_a t} \\
	e^{-i \omega_b t}\end{array} \right) \,, \quad 
\ket{\Psi_2} =\frac{1}{\sqrt{10}} \left( \begin{array}{c} e^{-i \omega_a t} \\
	-3e^{-i \omega_b t}\end{array} \right) \,, \\
\ket{\Psi_3} &=& \left( \begin{array}{c} 0 \\
	e^{-i \omega_b t}\end{array} \right) \,,
\end{eqnarray*}
%%
%
\begin{eqnarray*}
\rho &=& \sum_j P_j \ket{\Psi_j} \bra{\Psi_j} 	\\
	&=&\frac{3}{10} \ket{\Psi_1} \bra{\Psi_1} 
	+\frac{5}{10} \ket{\Psi_2} \bra{\Psi_2} 
	+\frac{2}{10}\ket{\Psi_3} \bra{\Psi_3} 	\\
&=& \frac{3}{20} \left( \begin{array}{cc} 1& e^{i \omega t} \\
	e^{-i \omega t} &1 \end{array} \right) 
	+\frac{1}{20} \left( \begin{array}{cc} 1& -3e^{i \omega t} \\
	-3e^{-i \omega t} &9 \end{array} \right) 	\\
	&& +\frac{4}{20} \left( \begin{array}{cc} 0& 0 \\
	0 &1 \end{array} \right) 	\\
	&=& \frac{1}{5} \left( \begin{array}{cc} 1& 0 \\
	0 &4 \end{array} \right) \,, 	\\
\rho^2 &=&\frac{1}{25} \left( \begin{array}{cc} 1& 0 \\
	0 &16 \end{array} \right) \neq \rho \,. 	
\end{eqnarray*}
%%
%
\begin{eqnarray*}
\bra{\Psi_1} \rho \ket{\Psi_1} &=& \frac{1}{10} 
	\left( e^{i\omega_a t} \,, e^{i\omega_b t} \right) 
	\left( \begin{array}{cc} 1& 0 \\ 0 &4 \end{array} \right) 
\left( \begin{array}{c} e^{i \omega_a t} \\ e^{-i \omega_b t} 
	\end{array} \right)  =\frac{1}{2} 	\,, 	\\
\bra{\Psi_2} \rho \ket{\Psi_2} &=& \frac{37}{50} \,, \qquad
	\bra{\Psi_3} \rho \ket{\Psi_3} =\frac{4}{5} \,.
\end{eqnarray*}
%%
Extra: La probabilidad de que el sistema este en el estado base es 
$\bra{\Psi_3} \rho \ket{\Psi_3} =4/5$ y de que este en el estado 
excitado es $(1\,, 0) \rho (1\,,0)^T =1/5$.
%\vspace{2cm}
%\end{widetext}

%\newpage
%%5%%%%%%%%60%%%%%%%%70%%%
\subsection{Principales Propiedades}
El operador de densidad obedece una ecuaci\'on: 
%
\begin{eqnarray}
\dot{\rho} &=& \sum_{\Psi} P_{\Psi} \left[ | \dot{\Psi} \rangle \bra{\Psi}  
	+ \ket{\Psi} \langle \dot{\Psi} | \right] 	\nonumber \\ 
&=& - \frac{i}{\hbar} \sum_{\Psi} P_{\Psi} \left[ 
	H \ketbra{\Psi}{\Psi} - \ketbra{\Psi}{\Psi} H \right] 
	\nonumber \\ 
&=& - \frac{i}{\hbar} (H \rho -\rho H) 	= - \frac{i}{\hbar} [H, \rho] 	\,.
\end{eqnarray}
%%
Esta es la ecuaci\'on de Schr\"odinger para el operador de densidad $\rho$, 
cuya soluci\'on es  
%
\begin{eqnarray}
\rho (t) = e^{-iHt/\hbar} \rho(0)  e^{iHt/\hbar} \,.
\end{eqnarray}
%%

Si se tiene una representaci\'on en t\'erminos de un conjunto completo de 
estados $\ket{n}$ cualquier operador $\hat{O}$ puede expanderse como 
%
\begin{eqnarray}
\hat{O} &=& \sum_{n,m} \bra{n} \hat{O} \ket{m} \ket{n} \bra{m}  
	= \sum_{n,m} O_{nm} \ket{n} \bra{m} \,,
\end{eqnarray}
%%
donde aplicamos la matriz identidad (matriz con unos en la diagonal 
principal y ceros fuera de ella) $\mathbf{1} =\sum_k \ket{k} \bra{k}$, y 
$O_{nm}$ son los elementos de matriz del operador $\hat{O}$. 

Consideremos, por ejemplo, el operador de momento dipolar 
$\mbd= e \mbr$: 
%
\begin{eqnarray}
e\mbr &=& e \sum_{n,m=1}^2 \bra{n} \mbr \ket{m} \ket{n} \bra{m}  
	\nonumber \\
	&=& e \left[ \bra{a} \mbr \ket{b} \ket{a} \bra{b} 
		+\bra{b} \mbr \ket{a} \ket{b} \bra{a}  \right] 		\nonumber \\
	&=& \mbd_{ab} \sigma_{ab} +\mbd_{ba} \sigma_{ba} \,,
\end{eqnarray}
%%
donde hicimos $\bra{a} \mbr \ket{a} = \bra{b} \mbr \ket{b} =0$, suponiendo 
que los estados at\'omicos poseen paridades opuestas. 

El valor esperado de un operador $\hat{O}$ esta dado por  
%
\begin{eqnarray}
\vxp{\hat{O}} &=& \sum_{\Psi} P_{\Psi} \bra{\Psi} \hat{O}\ket{\Psi}  \,,
\end{eqnarray}
%%
que toma en cuenta el peso $P_{\Psi}$ que cada estado tiene en el conjunto. 
Aplicamos tambi\'en  
$\rho = \sum_{\Psi} P_{\Psi} \ketbra{\Psi}{\Psi}$, 
%
\begin{eqnarray}
\vxp{\hat{O}} &=&  \sum_{\Psi} P_{\Psi} \sum_k \bra{\Psi} \hat{O}\ket{k}
		\braket{k}{\Psi} 	\nonumber \\ 
&=& \sum_k \sum_{\Psi} P_{\Psi} \braket{k}{\Psi}  \bra{\Psi} \hat{O}\ket{k} 
	= \sum_k \bra{k} \rho \hat{O} \ket{k} 	\nonumber  \\ 
&=& \sum_k ( \rho \hat{O} )_{kk} 	= \Tr ( \rho \hat{O} ) \,,
\end{eqnarray}
%%
donde la traza $\Tr$ indica la suma de los elementos en la diagonal principal 
de la matriz. 

Veamos la dependencia temporal del valor esperado del operador 
$\hat{O}$: 
%
\begin{eqnarray}
\vxp{\hat{O}(t)} &=& \Tr ( \rho(t) \hat{O}_S(0) ) 	
	=  \Tr ( e^{-iHt/\hbar} \rho(0)  e^{iHt/\hbar} \hat{O}_S ) \nonumber \\ 
	&=& \Tr (  \rho(0)  e^{iHt/\hbar} \hat{O}_S e^{-iHt/\hbar} ) \nonumber \\
	&=& \Tr ( \hat{O}_H(t) \rho(0) ) \,, 
\end{eqnarray}
%%
donde usamos la propiedad c\'iclica de la traza. Este resultado significa 
que podemos obtener los valores esperados ya sea en la imagen de 
Schr\"odinger (la primera forma) como en la imagen de Heisenberg.

Los elementos de matriz vienen dados por 
$\rho_{jk} = \bra{j} \rho \ket{k}$ y las ecuaciones para ellos son 
%
\begin{eqnarray}
\dot{\rho}_{nm} &=& \bra{n} \dot{\rho} \ket{m} 
	= - \frac{i}{\hbar}  \bra{n}  H \rho -\rho H \ket{m}		\nonumber \\ 
&=& - \frac{i}{\hbar} \sum_k \left[  
	\bra{n}  H \ketbra{k}{k} \rho \ket{m} 
	-  \bra{n} \rho \ketbra{k}{k}  H \ket{m}  \right] 	\nonumber \\ 
&=& - \frac{i}{\hbar} \sum_k \left[ H_{nk} \rho_{km} - \rho_{nk} H_{km} \right] \,.
\end{eqnarray}
%%

Para entender mejor la relaci\'on entre funci\'on de onda y operador de 
densidad y algunos de los resultados obtenidos para $\rho$ veamos el caso 
particular de un sistema de dos niveles con funci\'on de onda 
$\ket{\Psi} = C_a \ket{a} +C_b \ket{b}$.  Se tiene que 
%
\begin{eqnarray}
\rho &=& \left( C_a \ket{a} +C_b \ket{b} \right) 
	\left( C_a^\ast \bra{a} + \bra{b} C_b^\ast  \right)  	\nonumber \\ 
&=& C_a C_a^{\ast} \ketbra{a}{a} + C_a C_b^{\ast} \ketbra{a}{b} 
	+C_b C_a^{\ast} \ketbra{b}{a}  		\nonumber \\ 
	&& + C_b C_b^{\ast}\ketbra{b}{b} 		\nonumber \\ 
&=& \rho_{aa}  \ketbra{a}{a} +  \rho_{ab}  \ketbra{a}{b}  
	+  \rho_{ba} \ketbra{b}{a}  + \rho_{bb}  \ketbra{b}{b}  \,. \nonumber \\
\end{eqnarray}
%%

En notaci\'on matricial 
%
\begin{eqnarray}
\ket{a} &=&  \left( \begin{array}{c} 1 \\ 0 \end{array} \right) \,, \qquad 
	\ket{b} = \left( \begin{array}{c} 0 \\ 1 \end{array} \right) \,, \qquad 
	\ket{\Psi} = \left( \begin{array}{c} C_a \\ C_b \end{array} \right) \,. 
	\nonumber \\
\end{eqnarray}
%%
%
\begin{eqnarray}
\rho &=&  \left( \begin{array}{c} C_a \\ C_b \end{array} \right) 
	\left( \begin{array}{cc} C_a^{\ast} & C_b^{\ast} \end{array} \right) 	
	\nonumber \\ 
&=& \left( \begin{array}{cc}  C_a C_a^{\ast} & C_a C_b^{\ast}  \\ 
	C_b C_a^{\ast}  & C_b C_b^{\ast}   \end{array} \right) 	
	=  \left( \begin{array}{cc}  \rho_{aa} & \rho_{ab}  \\ 
	\rho_{ba}  & \rho_{bb}   \end{array} \right) 	\,,
\end{eqnarray}
%%
El t\'ermino $\ket{a}\bra{a}$, por mostrar un ejemplo, es 
%
\begin{eqnarray}
\ket{a}\bra{a} &=&  \left( \begin{array}{c} 1 \\ 0 \end{array} \right) 
	\left( \begin{array}{cc} 1 & 0 \end{array} \right) 	
	= \left( \begin{array}{cc}  1 & 0  \\ 0  & 0   \end{array} \right) \,.
\end{eqnarray}
%%
Este es el elemento de $\rho$ en la primera fila y primera columna. 
De manera similar se pueden identificar los otros elementos. 


Los t\'erminos de la diagonal, $\rho_{aa}$ y $\rho_{bb}$, son, 
respectivamente, las probabilidades de ocupar el estado superior e inferior 
del sistema, tambi\'en llamadas \textit{poblaciones}. Los t\'erminos fuera de 
la diagonal, $\rho_{ba} = \rho_{ab}^{\ast}$, son las llamadas 
\textit{coherencias}; en el caso de transiciones entre dos niveles 
representan los elementos de matriz del momento dipolar.  

Calculemos los valores esperados de los operadores elementales del sistema 
de dos niveles. Para ello utilizamos su representaci\'on de operadores de 
proyecci\'on, Secc~II.E: 
%
\begin{eqnarray}
\vxp{\sigma_+} &=& \Tr ( \rho \sigma_+ ) 	\nonumber\\ 
&=& \Tr \left[ \left( \rho_{aa}  \ketbra{a}{a} 
	+  \rho_{ab}  \ketbra{a}{b}  +  \rho_{ba} \ketbra{b}{a}  
	\right. \right. 	\nonumber\\  
  && \left. \left. + \rho_{bb}  \ketbra{b}{b}  \right) \ketbra{a}{b}   
	\right] 	\nonumber 	\\ 
&=& \Tr \left[  \rho_{aa} \ketbra{a}{b}  
	+  \rho_{ba}\ketbra{b}{b} \right] 	= \rho_{ba} \,, \\ 
\vxp{\sigma_-} &=& \vxp{\sigma_+}^\ast = \rho_{ab} \,.
\end{eqnarray}
%%
Con esto podemos definir la \textit{polarizaci\'on at\'omica}: 
%
\begin{eqnarray}
\vxp{\mbd} = \mbd_{ba} \rho_{ab} +\mbd_{ab} \rho_{ba} \,,
\end{eqnarray}
%%
y los componentes $x,y$ del operador de spin, 
%
\begin{eqnarray}
\vxp{\sigma_x} &=& \vxp{\sigma_+} + \vxp{\sigma_-} 
	= \rho_{ab} + \rho_{ba} 	\,,	\\ 
\vxp{\sigma_y}	&=& i \vxp{\sigma_+} -i \vxp{\sigma_-} 
	= i ( \rho_{ab} - \rho_{ba} ) 	\,.
\end{eqnarray}
%%

Con la componente $z$ definimos la \textit{inversi\'on at\'omica}: 
%
\begin{eqnarray}
\vxp{\sigma_z} &=& \Tr ( \rho \sigma_z ) 	\nonumber\\ 
&=& \Tr \left[ \left( \rho_{aa} \ketbra{a}{a} 
	+  \rho_{ab}  \ketbra{a}{b}  +  \rho_{ba} \ketbra{b}{a} 
	\right. \right. 	\nonumber\\  
  && \left. \left. + \rho_{bb} \ketbra{b}{b}  \right) 
  \left( \ketbra{a}{a} - \ketbra{b}{b} \right) 
	\right] 	\nonumber 	\\ 
&=& \Tr \left[ \rho_{aa} \ketbra{a}{a} 
	-  \rho_{ab} \ketbra{a}{b} + \rho_{ba} \ketbra{b}{a} \right. \nonumber \\
	 && \left. - \rho_{bb} \ketbra{b}{b}  \right] 	\nonumber \\ 
&=& \rho_{aa} - \rho_{bb} 	\,.
\end{eqnarray}
%%
 
Partiendo del modelo dado por las Ecs.(\ref{eq:decayAmpl}), las ecuaciones 
para los elementos del operador de densidad son:
%
\begin{eqnarray}
\dot{\rho}_{aa} &=& \dot{C}_a C_a^{\ast} + C_a \dot{C}_a^{\ast}   
    = -\gamma_a \rho_{aa} + \frac{i}{2} \Omega_0 (\rho_{ba} - \rho_{ab}) \,, 
    	\nonumber \\ 
\dot{\rho}_{bb} &=& \dot{C}_b C_b^{\ast} + C_b \dot{C}_b^{\ast}  
  = -\gamma_b \rho_{bb} - \frac{i}{2} \Omega_0 (\rho_{ba} - \rho_{ab}) \,, 
  	\nonumber \\ 
\dot{\rho}_{ab} &=& \dot{C}_a C_b^{\ast} + C_a \dot{C}_b^{\ast}  \nonumber \\ 
  	&=& - (\gamma_{ab} +i\Delta) \rho_{ab} 
  	- \frac{i}{2} \Omega_0 (\rho_{aa} - \rho_{bb})  \,.
\end{eqnarray}
%%

%%5%%%%%%%%60%%%%%%%%70%%%
\subsection{Emisi\'on Espont\'anea} 
En la Secci\'on II-B vimos el caso de excitaci\'on de un nivel hacia un cont\'inuo. 
El acoplamiento de uno o mas estados con un conjunto grande de estados muy 
cercanos llev\'o a que la probabilidad de permanecer en un estado decae en el 
tiempo de forma exponencial. En la Secci\'on III-A a\~nadimos la dispaci\'on de 
manera fenomenol\'ogica, es decir, sabemos que la hay pero no hurgamos en 
los detalles de su origen.   

La emisi\'on espont\'anea es el proceso por el cual un \'atomo en estado 
excitado decae a un estado de menor energ\'ia emitiendo un fot\'on. La 
cuesti\'on es: ?`c\'omo ocurre este decaimiento? Este es un tema avanzado 
de la \'optica cu\'antica. Sin embargo, la clave est\'a en el acoplamiento de un 
sistema peque\~no (el \'atomo) con uno muy grande, un entorno. En este 
caso, el sistema grande consta de los modos del campo electromagn\'etico 
cuantizado de una cavidad tan grande que contiene un n\'umero muy grande y 
cuasi-cont\'inuo de esos modos. En la pr\'actica, no nos interesa la evoluci\'on 
de esos modos, s\'olo la del \'atomo. Por un proceso de \textit{eliminaci\'on} de 
los modos del entorno la din\'amica se reduce a la del \'atomo con un precio a 
pagar: la probabilidad de que el \'atomo permanezca en el estado excitado 
decae como $P_a \propto e^{-\gamma t}$, donde 
$\gamma \propto \wp^2 \omega_{ab}^3$, $\wp$ es el momento dipolar y 
$\omega_{ab}$ es la frecuencia de la transici\'on entre los niveles inicial y 
final.  Recu\'erdese que esta $\gamma$ es la principal dificultad para realizar 
l\'aseres en UV y X con tecnolog\'ia similar a los del visible, porque la vida 
media del estado excitado es tan corta que no alcanza a sostener la 
poblaci\'on necesaria para la inversi\'on.   

El mencionado proceso de eliminaci\'on del entorno, usando el m\'etodo 
de Weisskopf y Wigner que considera acoplamiento d\'ebil entre \'atomo 
y modos del entorno y en la aproximaci\'on de Markov, lleva al mencionado 
decaimiento exponencial. El m\'etodo Weisskopf-Wigner en el contexto 
del operador de densidad lleva a agregar a la ecuaci\'on para el operador 
de densidad del \'atomo los t\'erminos 
%
\begin{eqnarray} 	\label{eq:mEq2LA_bb}
\mathcal{L}_a \rho &=& \frac{\gamma_a}{2} (\bar{n} +1) 
	\left( 2 \sigma_- \rho \sigma_+ -\sigma_+ \sigma_- \rho 
	-\rho \sigma_+ \sigma_- \right) 	\nonumber \\
	&& + \frac{\gamma_a}{2} \bar{n}  
	\left( 2 \sigma_+ \rho \sigma_- -\sigma_- \sigma_+ \rho 
	-\rho \sigma_- \sigma_+ \right)\,, 
\end{eqnarray}
%%
donde $\bar{n}$ es el n\'umero promedio de fotones a temperatura 
$T$. Si, adem\'as, el entorno tiene temperatura cero, $\bar{n} =0$ lo 
anterior se reduce a 
%
\begin{eqnarray}
\mathcal{L}_a \rho = \frac{\gamma_a}{2} \left( 2 \sigma_- \rho \sigma_+ 
	-\sigma_+ \sigma_- \rho -\rho \sigma_+ \sigma_- \right) \,, 
\end{eqnarray}
%%
que es una excelente aproximaci\'on a frecuencias \'opticas a\'un a 
temperatura ambiente, $T = 300$K. Este t\'ermino justifica plenamente 
el decaimiento por emisi\'on espont\'anea que inicialmente se a\~nadi\'o 
fenomenol\'ogicamente. 

\subsubsection{Ecuaciones de Bloch} 
Una vez que hemos introducido formalmente la emisi\'on espont\'anea a 
nuestro modelo de interacci\'on \'atomo-l\'aser procedemos a obtener las 
ecuaciones de movimiento para los elementos de la matriz de densidad. 
Estas son llamadas Ecuaciones de Bloch, que generalizan las ecuaciones 
de un esp\'in acoplado a un campo magn\'etico. Las ecuaciones de Bloch 
provienen de la llamada \textit{ecuaci\'on maestra} 
%
\begin{eqnarray}
\dot{\rho} = -\frac{i}{\hbar} [H,\rho]  +\frac{\gamma_a}{2} 
	\left( 2 \sigma_- \rho \sigma_+ -\sigma_+ \sigma_- \rho 
	-\rho \sigma_+ \sigma_- \right) \,. 	\nonumber \\
\end{eqnarray}
%%
Nos falta solamente especificar el Hamiltoniano (Secc. II.E), 
%
\begin{eqnarray}
 H = \hbar \omega_0 \sigma_{aa} + \frac{\Omega_0}{2} 
	(\sigma_+ e^{-i\nu t} 	+\sigma_- e^{i\nu t}) \,,
\end{eqnarray}
%%
donde el primer t\'ermino es la energ\'ia del \'atomo no perturbado (con 
energ\'ia de transici\'on $\hbar \omega_0$) y el segundo es la energ\'ia de 
interacci\'on del \'atomo con el l\'aser. N\'otese que desplazamos el origen 
de la energ\'ia del \'atomo de la mitad de la transici\'on a la del estado base.  

De aqu\'i, proyectando por $\bra{j}$ por la izquierda y por $\ket{k}$ por la 
derecha, donde $j,k =a,b$, en ambos lados de la ecuaci\'on, obtenemos las ecuaciones para $\rho_{jk} = \bra{j} \rho \ket{k}$: 
\newline \textcolor{red}{Pendiente}

La presencia de las frecuencias r\'apidas $\omega_0$ y $\nu$ hace 
inconveniente usar estas ecuaciones para soluci\'on num\'erica. Entonces, 
si definimos $\tilde{\rho}_{ab}(t) = \rho_{ab} e^{i\nu t}$ y 
$\Delta = \omega_0 -\nu$, obtenemos las ecuaciones
%
\begin{eqnarray} 	\label{eq:BlochEq_rho}
\dot{\tilde{\rho}}_{ab} 
	&=& - \left( \frac{\gamma_a}{2} +i\Delta \right) \tilde{\rho}_{ab}
	-i\frac{\Omega_0}{2} (\rho_{aa} -\rho_{bb}) 	\,, \nonumber \\ 
\dot{\tilde{\rho}}_{aa} 
	&=& -\gamma_a \rho_{aa} + i\frac{\Omega_0}{2} \tilde{\rho}_{ba} 
	- i\frac{\Omega_0}{2} \tilde{\rho}_{ab} 	\,, 	\nonumber \\ 
\dot{\tilde{\rho}}_{bb} 
	&=& \gamma_a \rho_{aa} - i\frac{\Omega_0}{2} \tilde{\rho}_{ba} 
	+ i\frac{\Omega_0}{2} \tilde{\rho}_{ab} 	\,, 	\\ 
\dot{\tilde{\rho}}_{ab} 
	&=& 	 \dot{\tilde{\rho}}_{ba}^\ast  \,. \nonumber
\end{eqnarray}
%%
Decimos que estas ecuaciones evolucionan en un marco de referencia 
rotando a la frecuencia del l\'aser $\nu$. \textcolor{blue}{Esto es similar a 
fijarnos \'unicamente en el movimiento lento de un trompo ignorando su 
rotaci\'on r\'apida, o solamente la inclinaci\'on de la Tierra en su \'orbita 
alrededor del Sol, ignorando su rotaci\'on diaria.  }

Posteriormente a\~nadiremos el efecto de temperatura a la ecuaci\'on 
para $\rho$. 

Las Ecs.(\ref{eq:BlochEq_rho}) forman un conjunto de 4 ecuaciones 
homog\'eneas. Para resolverlo conviene reducirlo eliminando la variable 
$\rho_{bb}$  usando la conservaci\'on de la probabilidad 
$\rho_{aa} +\rho_{bb} =1$. [La forma obtenida no es muy com\'un en la 
literatura pero adjuntamos la soluci\'on en una forma alternativa 
presentada abajo.] Las nuevas ecuaciones son: 
%
\begin{eqnarray} 	\label{eq:BlochEq_rhoV2}
\dot{\tilde{\rho}}_{ab} 
	&=& - \left( \frac{\gamma_a}{2} +i\Delta \right) \tilde{\rho}_{ab}
	-i\frac{\Omega_0}{2} (2\rho_{aa} -1) 	\,, \nonumber \\ 
\dot{\tilde{\rho}}_{ba} 
	&=& - \left( \frac{\gamma_a}{2} -i\Delta \right) \tilde{\rho}_{ba}
	+i\frac{\Omega_0}{2} (2\rho_{aa} -1) 	\,, \nonumber \\ 
\dot{\tilde{\rho}}_{aa} 
	&=& -\gamma_a \rho_{aa} + i\frac{\Omega_0}{2} \tilde{\rho}_{ba} 
	- i\frac{\Omega_0}{2} \tilde{\rho}_{ab} 	\,, 	
\end{eqnarray}
%%

\subsubsection{Soluciones en Estado Estacionario} 
Los elementos $\tilde{\rho}_{jk}(t)$ evolucionan en el tiempo en funci\'on de los 
par\'ametros $\gamma_a$, $\Omega_0$ y $\Delta$. Obtenemos, sin embargo, excelente informaci\'on del sistema con las soluciones en estado estacionario, 
es decir, a tiempos muy largos, cuando $t \to \infty$. En este caso los 
elementos de matriz llegan a un valor constante 
$\tilde{\rho}_{jk}(t \to \infty) = \tilde{\rho}_{jk}^{st}$, donde 
$\dot{\tilde{\rho}}_{jk} (t \to \infty) =0$. Resolviendo el sistema de ecuaciones, 
%
\begin{subequations}
\begin{eqnarray}
\tilde{\rho}_{ab}^{st} &=& \frac{-\Omega_0 [2\Delta + i \gamma]}
	{2\Omega_0^2 +\gamma^2 +4\Delta^2} \,, 
	\qquad \tilde{\rho}_{ba}^{st} = \left( \tilde{\rho}_{ab}^{st} \right)^\ast , 
	\quad \\ 
\tilde{\rho}_{bb}^{st} &=& \frac{\Omega_0^2 +\gamma^2 +4\Delta^2}
	{2\Omega_0^2 +\gamma^2 +4\Delta^2}  \,, \\ 
\tilde{\rho}_{aa}^{st} &=& \frac{\Omega_0^2}{2\Omega_0^2 
	+\gamma^2 +4\Delta^2}	\,.
\end{eqnarray}
\end{subequations}
%%

Veamos algunas propiedades y casos de particular inter\'es. 
\newline (i) Puesto que el estado $\ket{a}$ decae al estado $\ket{b}$ la 
poblaci\'on total se conserva a todo tiempo, $\rho_{aa} +\rho_{bb} = 1$, 
entonces $\dot{\rho}_{aa} +\dot{\rho}_{bb} = 0$. 
\newline (ii) $\rho_{aa}$ est\'a relacionado con la inversi\'on de la poblaci\'on 
por la acci\'on del l\'aser. La curva es de forma Lorentziana, ensanchada por 
la intensidad $\Omega_0^2 \propto E_0^2$ del l\'aser. A este efecto se le 
denomina \textit{ensanchamiento por potencia}. 
\newline (iii) Cuando $\Omega_0 \ll \gamma$ el \'atomo casi no absorbe 
energ\'ia, por lo que permanece en su estado base, $\rho_{aa}^{st} \to 0$, 
$\rho_{bb}^{st} \simeq 1$. Una situaci\'on similar ocurre cuando el l\'aser 
est\'a muy fuera de sinton\'ia con la transici\'on at\'omica, 
$\Delta \gg \Omega_0$. 
\newline (iv) Cuando $\Omega_0 \gg \gamma, \Delta$ la absorci\'on es alta. 
Sin embargo, en estado estacionario, la probabilidad de ocupaci\'on del 
estado superior siempre es menor o igual a la del estado base, obteniendo 
la igualdad cuando $\Delta =0$, $\rho_{aa}^{st} = \rho_{bb}^{st} = 1/2$. 
\newline (v) $\rho_{ab}$ es la coherencia  o polarizaci\'on del \'atomo (versi\'on microsc\'opica del \'indice de refracci\'on complejo de un medio). La parte real 
es llamada \textit{dispersi\'on} y la imaginaria es la \textit{absorci\'on}. 
$\rho_{ab}$ es puramente imaginaria cuando $\Delta =0$. 
\newline (vi) Cuando $\Omega \sim \gamma$ la interacci\'on \'atomo-l\'aser 
es suficientemente fuerte, entrando en un r\'egimen no lineal llamado de \textit{saturaci\'on}, donde $\rho_{aa}$ ya no es despreciable. En este 
r\'egimen es \'util el concepto de \textit{inversi\'on} de poblaci\'on, que no 
aparece en el modelo cl\'asico del electr\'on, v\'alido en el r\'egimen 
$\Omega_0 \ll \gamma_a$.  \newline (vii) Finalmente, es conveniente conocer 
valores t\'ipicos de los par\'ametros $\gamma_a$, $\Omega_0$ y $\Delta$ en experimentos de \'optica cu\'antica. Usualmente el par\'ametro de referencia 
es la raz\'on de emisi\'on espont\'anea $\gamma_a$, que para transiciones 
\'opticas es del orden de MHz. La frecuencia de Rabi $\Omega_0$ puede 
ser hasta unas veces $\gamma_a$. Si es m\'as fuerte tal vez sea necesario 
hacer ajustes a nuestro modelo \'atomo-l\'aser. La desinton\'ia tambi\'en 
puede ser unas veces $\gamma_a$.  

\subsubsection{Soluciones dependientes del tiempo} 
Las ecuaciones de Bloch son suficientemente complicadas para obtener 
soluciones para valores arbitrarios de los par\'ametros 
$\gamma_a,\Omega_0,\Delta$. Un caso particularemente com\'un y \'util 
es cuando el l\'aser est\'a en resonancia con el \'atomo, $\Delta =0$. Una 
forma de resolver las ecuaciones de Bloch es mediante transformadas de 
Laplace, con condiciones iniciales $\tilde{\rho}_{bb} (0) =1$ y 
$\tilde{\rho}_{aa} (0) = \tilde{\rho}_{ab} (0) = \tilde{\rho}_{ba} (0) = 0$. 
Obtenemos
%
\begin{subequations} 	%\label{eq:BlochSolsApp}
\begin{eqnarray} 	
\tilde{\rho}_{ab}(t) &=& - i \frac{ Y/ \sqrt{2} }{1+Y^2}  f(t)
  \mp i \sqrt{2} Y \frac{\gamma_a}{4 \delta} e^{-3\gamma_a t/4} \sinh{\delta t}	
   \,,  	\nonumber \\   \\
%
\tilde{\rho}_{aa}(t) &=& \frac{ Y^2/2 }{1+ Y^2}  f(t)  \,, 	\\ 
%
\tilde{\rho}_{bb}(t) &=& 1 -\frac{ Y^2/2 }{1+ Y^2}   f(t)  \,, 
\end{eqnarray}
%%
donde
%
\begin{eqnarray}
f(t) &=& 1 -e^{-3\gamma_a t/4} \left( \cosh{\delta t}  
	+ \frac{3 \gamma_a}{4 \delta} \sinh{\delta t} \right) \,, \nonumber \\ \\
Y &=& \frac{ \sqrt{2} \Omega_0 }{\gamma_a} 	\,, \qquad 
	\delta  = \frac{\gamma_a}{4} \sqrt{1 - 8Y^2} \,. 
\end{eqnarray}
\end{subequations} 
%% 

Cuando la excitaci\'on es d\'ebil, $ 8Y^2 \leq 1$, o bien 
$\Omega_0 < \gamma /2$ las soluciones var\'ian mon\'otonamente. 
Cuando $ 8Y^2 \geq 1$, o bien $\Omega_0 > \gamma /2$, se tiene 
$\cosh{\delta \tau} \to \cos{\delta' \tau}$ y 
$\sinh{\delta \tau}/\delta \to \sin{\delta' \tau}/\delta'$, donde 
$\delta = i \delta' = (\gamma_a /\sqrt{8Y^2 -1}$. Las $\tilde{\rho}_{jk}(t)$ 
se hacen oscilatorias para $\Omega_0 \gg \gamma_a$. 

Otro caso en el que es posible una soluci\'on anal\'itica es cuando 
$\gamma_a \to 0$ \ldots 

%%5%%%%%%%%60%%%%%%%%70%%%
\subsection{Forma Alternativa de las Ecuaciones de Bloch 1}  
En la pr\'actica, hay otras formas de representar las ecuaciones de Bloch. 
El primer paso es recordar que $\vxp{\tilde{\sigma}_{jk}} = \tilde{\rho}_{kj}$ 
(para las poblaciones podemos omitir la tilde). Luego definimos la llamada 
\textit{inversi\'on} de poblaci\'on: 
$\vxp{\sigma_z} = \vxp{\sigma_{aa}} -\vxp{\sigma_{bb}}$. Junto con la 
conservaci\'on de la probabilidad $1 = \vxp{\sigma_{aa}} +\vxp{\sigma_{bb}}$ 
relacionamos poblaciones e inversi\'on como 
$\vxp{\sigma_{aa}} = (1+ \vxp{\sigma_z})/2$ y 
$\vxp{\sigma_{bb}} = (1- \vxp{\sigma_z})/2$. De esta manera, el sistema 
(\ref{eq:BlochEq_rho}) se reduce a tres ecuaciones inhomog\'eneas: 
%
\begin{eqnarray} 	\label{eq:BlochEq_+-z}
\vxp{\dot{\sigma}_z} &=& - \gamma_a ( \vxp{\sigma_z} +1)  
	+ i \Omega_0 \vxp{\tilde{\sigma}_+} 
	- i \Omega_0 \vxp{\tilde{\sigma}_-}	\,, \nonumber \\ 
\vxp{\dot{\tilde{\sigma}}_-}
	&=& - \left( \frac{\gamma_a}{2} +i\Delta \right) \vxp{\tilde{\sigma}_-} 
	-i\frac{\Omega_0}{2} \vxp{\sigma_z} 	\,, \nonumber \\ 
\vxp{\dot{\tilde{\sigma}}_+}
	&=& - \left( \frac{\gamma_a}{2} -i\Delta \right) 
	\vxp{\tilde{\sigma}_+} +i\frac{\Omega_0}{2} \vxp{\sigma_z} 	\,, 
\end{eqnarray}
%%
Adjuntamos la soluci\'on en nota aparte.  

%%5%%%%%%%%60%%%%%%%%70%%%
\subsection{Forma Alternativa de las Ecuaciones de Bloch: 
Modelo Vectorial} 
Una descripci\'on muy \'util e ilustrativa de la interacci\'on de un sistema 
de dos niveles con un campo es el modelo vectorial (Feynman, Vernon, 
Hellwart??). Primero introducimos las variables 
%
\begin{subequations}
\begin{eqnarray}
\tilde{u} &=& \vxp{\tilde{\sigma}_x} = \tilde{\rho}_{ab} + \tilde{\rho}_{ba} \,, \\ 
\tilde{v} &=& \vxp{\tilde{\sigma}_y} = i \tilde{\rho}_{ab} - i \tilde{\rho}_{ba} \,, \\ 
\tilde{w} &=& \vxp{\sigma_z} = \tilde{\rho}_{aa} + \tilde{\rho}_{bb} \,, 
\end{eqnarray}
\end{subequations}
%%
con lo que definimos el \textit{vector de Bloch} 
%
\begin{eqnarray}
\mbB &=& \tilde{u} \mbe_x + \tilde{v} \mbe_y +  \tilde{w} \mbe_z 
	= \Tr [ \tilde{\rho} \tilde{\sigma} ] \,.
\end{eqnarray}
%%

Con las ecuaciones (\ref{eq:BlochEq_rho}) formamos el conjunto de 
ecuaciones para los componentes de $\mbB$: 
%
\begin{subequations}
\begin{eqnarray}
\dot{\tilde{u}} &=& -\frac{\gamma_a}{2} \tilde{u} - \Delta \tilde{v} \,, \\ 
\dot{\tilde{v}} &=& \Delta \tilde{u} - \frac{\gamma_a}{2} \tilde{v} 
	+\Omega_0 \tilde{w} \,, \\
\dot{\tilde{w}} &=& -\gamma_a (\tilde{w} +1) -\Omega_0 \tilde{v} \,. 
\end{eqnarray}
\end{subequations}
%%

Las soluciones en estado estacionario son
%
\begin{subequations}
\begin{eqnarray}
\tilde{u}_{st} &=& \frac{\Delta \Omega_0}
	{\Omega_0^2/2 +\Delta^2 +(\gamma_a/2)^2 } \,, \\ 
\tilde{v}_{st} &=& -\frac{ \Omega_0 \gamma_a/2}
	{\Omega_0^2/2 +\Delta^2 +(\gamma_a/2)^2 } \,, \\
\tilde{w}_{st} &=& - \frac{(\gamma_a/2)^2 +\Delta^2}
	{\Omega_0^2/2 +\Delta^2 +(\gamma_a/2)^2 }\,, 
\end{eqnarray}
\end{subequations}
%%

Hay que notar que los componentes del vector de Bloch son variables 
reales, lo que facilita la interpretaci\'on del sistema como un vector en 
un espacio real en 3D.

Varios aspectos de estas soluciones son muy importantes que conviene 
tener siempre muy presentes. (i) El valor m\'aximo de la inversi\'on de 
poblaci\'on $\tilde{w}_{st}$ se obtiene cuando $\Delta=0$ y 
$\Omega_0 \gg \gamma_a$, dando $\tilde{w}_st =0$, o bien 
$\rho_{aa} = \rho_{bb} =1/2$; (ii) cuando $\Delta=0$ el movimiento del 
vector de Bloch ocurre en el plano $y,z$ ($\tilde{u}(t)=0$); (iii)la inversi\'on 
$\tilde{w}$ no tiene contraparte cl\'asica (donde el \'atomo no tiene niveles 
de energ\'ia); (iv) las coherencias $\rho_{ab}$ y $\rho_{ba}$ representan 
la polarizaci\'on inducida en el medio, 
$\tilde{\rho}_{ab} = (\tilde{u} -i \tilde{v})/2$.

El modelo vectorial permite visualizar la no linealidad de la interacci\'on  
\'atomo-l\'aser utilizando el concepto de \textit{saturaci\'on}. Esta indica 
que un cierto nivel de intensidad del l\'aser es suficiente para inducir 
una respuesta no lineal en el \'atomo. Definiremos el par\'ametro de 
saturaci\'on $I_S$ tal que 
%
\begin{eqnarray*}
\tilde{w}_{st} &=& - \left[ 1 + 
	\frac{\Omega_0^2 /2}{(\gamma_a/2)^2 +\Delta^2)} \right]^{-1} \\ 
&=& - \left[ 1 + 
	\frac{\Omega_0^2 /2}{(\gamma_a/2)^2} \cdot 
	\frac{1}{1+ \frac{4\Delta^2}{\gamma_a^2}} \right]^{-1} \,.
\end{eqnarray*}
%%
Definiendo  
%
\begin{eqnarray}
I_S = \frac{2\Omega_0^2}{\gamma_a^2} \,, 	\qquad 
\mcL (\Delta) = \left[ 1+ \frac{4\Delta^2}{\gamma_a^2} \right]^{-1} \,, 
\end{eqnarray}
%%
entonces 
%
\begin{subequations}
\begin{eqnarray}
\tilde{w}_{st} &=& \frac{-1}{ 1+ I_S \mcL(\Delta) } \,,  \\ 
\tilde{v}_{st} &=& -\frac{ 2(\Omega_0/\gamma_a) \mcL(\Delta)} 
	{1+ I_S \mcL(\Delta) } \,, \\
\tilde{u}_{st} &=& \frac{ (4\Delta \Omega_0 /\gamma_a) \mcL(\Delta) }
	{1+ I_S \mcL(\Delta) } \,, 
\end{eqnarray}
\end{subequations}
%%
El par\'ametro de saturaci\'on $I_S$ es una medida adimensional de la 
intensidad requerida para saturar el \'atomo, es decir, cuando su capacidad 
de absorci\'on disminuye. Esto empieza a ocurrir cuando $I_S =1$ 
($\Omega_0 = \gamma_a/\sqrt{2}$).

\textcolor{red}{\'Area de un pulso:} 

%%5%%%%%%%%60%%%%%%%%70%%%
\subsection{Desfasamiento, Eliminaci\'on Adiab\'atica y Ecuaciones 
de Raz\'on} 
Otro proceso de gran importancia es el de las colisiones el\'asticas entre 
los \'atomos, las cuales solamente producen un desfasamiento entre 
diferentes emisores, sin afectar los procesos de absorci\'on y emisi\'on. 
Ahora debemos a\~nadir el t\'ermino 
%
\begin{eqnarray}
\mathcal{L}_p \tilde{\rho} = 
	 (\gamma_p /2) (\sigma_z \tilde{\rho} \sigma_z -\tilde{\rho}) 
\end{eqnarray}
%%
a la ecuaci\'on para $\tilde{\rho}$, donde $\gamma_p$ se interpreta como 
la frecuencia de las colisiones, que se refleja \'unicamente en las 
ecuaciones para las coherencias a\~nadiendo t\'erminos de la forma 
$(\dot{{\tilde{\rho}}}_{ab})_p = -\gamma_p \tilde{\rho}_{ab}$ 
[$(\dot{{\tilde{\rho}}}_{aa})_p = (\dot{{\tilde{\rho}}}_{bb})_p = 0$]. 
Por lo tanto, las coherencias tienen decaimiento adicional de la forma 
$\tilde{\rho}_{ab} (t) \propto \tilde{\rho}_{ab} e^{- \gamma t}$), donde  
%
\begin{eqnarray}
\gamma = \frac{\gamma_a}{2} +\gamma_p  \,. 
\end{eqnarray}
%%
Las ecuaciones son ahora, 
%
\begin{eqnarray} 	%\label{eq:BlochEq_rho}
\dot{\tilde{\rho}}_{ab} 
	&=& - \left( \gamma +i\Delta \right) \tilde{\rho}_{ab}
	-i\frac{\Omega_0}{2} (\rho_{aa} -\rho_{bb}) 	\,, \nonumber \\ 
\dot{\tilde{\rho}}_{aa} 
	&=& -\gamma_a \rho_{aa} + i\frac{\Omega_0}{2} \tilde{\rho}_{ba} 
	- i\frac{\Omega_0}{2} \tilde{\rho}_{ab} 	\,, 	\nonumber \\ 
\dot{\tilde{\rho}}_{bb} 
	&=& \gamma_a \rho_{aa} - i\frac{\Omega_0}{2} \tilde{\rho}_{ba} 
	+ i\frac{\Omega_0}{2} \tilde{\rho}_{ab} 	\,,
\end{eqnarray}
%%

Si $\gamma_p \gg \gamma_a$, es decir, el desfasamiento es mucho 
m\'as r\'apido que la emisi\'on espont\'anea, es posible simplificar el 
sistema eliminando la variable r\'apida $\tilde{\rho}_{ab}$, que adquiere 
su estado estacionario m\'as r\'apidamente que las poblaciones. 
Haciendo $\dot{\tilde{\rho}}_{ab} =0$, entendiendo que 
$\tilde{\rho}_{ab} \to \tilde{\rho}_{ab}^{st}$; 
entonces 
%
\begin{eqnarray}
\tilde{\rho}_{ab} &=& 
	- i \frac{\Omega_0/2}{\gamma +i\Delta} (\rho_{aa} - \rho_{bb}) \,, 
\end{eqnarray}
%%
que sustituimos en las ecuaciones para las poblaciones: 
%
\begin{subequations}
\begin{eqnarray}
\dot{\rho}_{aa} &=&  -\gamma_a \rho_{aa}  
	- \frac{\Omega_0^2}{4} ( \rho_{aa} - \rho_{bb})  
	\left[ \frac{1}{\gamma -i \Delta} + \frac{1}{\gamma +i\Delta} \right]   
	\nonumber \\ 
&=& -\gamma_a \rho_{aa} 
	- \frac{ \gamma \Omega_0^2/2 }{ \Delta^2 +\gamma^2 } 
	( \rho_{aa} - \rho_{bb}) 	\,, 
\end{eqnarray}
%%
%
\begin{eqnarray}
\dot{\rho}_{bb} &=&  \gamma_a \rho_{aa}  
	+ \frac{\Omega_0^2}{4} ( \rho_{aa} - \rho_{bb})  
	\left[ \frac{1}{\gamma -i \Delta} + \frac{1}{\gamma +i\Delta} \right]    
	\nonumber \\ 
&=& \gamma_a \rho_{aa} 
	+ \frac{ \gamma \Omega_0^2/2 }{ \Delta^2 +\gamma^2 } 
	( \rho_{aa} - \rho_{bb}) 	\,. 
\end{eqnarray}
\end{subequations}
%%
Estas son llamadas \textit{Ecuaciones de Raz\'on}, en las que solamente 
encontramos a las poblaciones o probabilidades de ocupaci\'on de los 
niveles, debido a que las colisiones destruyen la coherencia entre ellos. 

Este procedimiento, llamado \textit{eliminaci\'on adiab\'atica} es muy \'util 
siempre que un sistema f\'isico presenta dos o mas escalas de tiempo 
muy diferentes entre s\'i, es decir, una es mucho mas r\'apida que la otra. 
En estos casos se dice que la variable r\'apida (eliminada) sigue 
adiab\'aticamente a la variable lenta. 

En s\'olidos hay procesos an\'alogos a las colisiones, como son las 
vibraciones de la red, defectos, etc. Si $\gamma_p$ es similar o mayor a 
$\gamma$ podemos aplicar el proceso de eliminaci\'on adiab\'atica visto 
arriba. 

\textbf{Ejemplo. Fase Relativa Aleatoria.}
Una forma un poco m\'as formal de ver el efecto de desfasamiento es el  
siguiente. Con frecuencia hay que realizar promedios estad\'isticos a las 
variables de inter\'es adem\'as de los promedios cu\'anticos. Considere 
un conjunto de vectores de estado de dos niveles $\ket{\psi}_j$ todos 
con las mismas probabilidades $|C_a|^2$ y $|C_b|^2$, pero que tienen 
una fase relativa aleatoria $\phi_j$ entre las amplitudes de probabilidad: 
$\ket{\phi}_j = c_a \ket{a} + e^{i \phi_j} c_b \ket{b}$. 

Para un sistema de $N$ \'atomos la probabilidad $P_j$ del $j$-\'esimo estado 
es $1/N$, dado que todos las fases $\phi_j$ son igualmente probables. El 
elemento no diagonal $\rho_{ab}$ para este problema es 
%
\begin{eqnarray*}
\rho_{ab} &=& \sum_j P_j c_{aj} c_{bj}^\ast 
	= c_a c_b^\ast \frac{1}{N} \sum_{j=1}^N e^{-i \phi_j} 	\\ 
	&=& c_a c_b^\ast \frac{1}{N}  \int_0^{2\pi} d\phi \, e^{-i \phi} = 0 \,,
\end{eqnarray*}
%%
es decir, la coherencia at\'omica entre los niveles $a$ y $b$ desaparece 
completamente en el proceso de promediar. 

%%5%%%%%%%%60%%%%%%%%70%%%
\subsection{Excitaci\'on T\'ermica y/o Incoherente} 
Anteriormente vimos la Ec.~(\ref{eq:mEq2LA_bb}) para la interacci\'on 
de un \'atomo de dos niveles con un campo EM coherente y un campo a  
temperatura $T$, que aparece mediante el n\'umero de fotones t\'ermicos 
$\bar{n}$. Supongamos ahora que el \'atomo es excitado solamente por 
un campo incoherente, por lo que descartamos el t\'ermino 
$(-i/\hbar)[H,\rho]$ y el campo, en lugar de deberse a la temperatura 
suponemos que se debe a un l\'aser muy ruidoso. Hacemos, entonces, 
el remplazo $\gamma_a \bar{n}/2 \to \Gamma$, y a\~nadimos a la 
ecuaci\'on maestra el t\'ermino
%
\begin{eqnarray} 	\label{eq:mEq2LA_inc}
\mathcal{L}_{\Gamma} \rho &=& \Gamma   \left( 2 \sigma_+ \rho \sigma_- 
	-\sigma_- \sigma_+ \rho -\rho \sigma_- \sigma_+ \right)\,, 
\end{eqnarray}
%%
Obtener ecuaciones y soluci\'on \ldots

%%%%%%%%10%%%%%%%%20%%%%%%%%30%%%%%%%%40%%%%%%%%50%%%%%%%%60%%%%%%%%70%%%
\section{Coherencia vs Incoherencia en el \'Atomo de Dos NIveles} 
Uno de los grandes objetivos de la \'optica cu\'antica es determinar las 
fluctuaciones de la luz y de las fuentes que la producen. Para describir el 
papel que tienen las fluctuaciones separamos los operadores at\'omicos 
en una parte dada por su valor esperado (su media) a tiempos largos, 
$ \vxp{\sigma_{jk}}_{st} =\alpha_{jk}$, mas 
las fluctuaciones, $\Delta \sigma_{jk}(t)$, es decir,  
%
\begin{eqnarray} 	\label{eq:dipoleSplit}
\sigma_{jk} (t) =\alpha_{jk}  +\Delta \sigma_{jk} (t) 	\,, 
\end{eqnarray}
%%
donde, por definici\'on, $\vxp{\Delta \sigma_{jk}(t)} =0$,  tal que 
%
\begin{eqnarray}
\vxp{\sigma_{eg} (0) \sigma_{ge} (\tau)}_{st}  
&=& |\alpha_{eg} |^2 
	+\vxp{\Delta \sigma_{eg} (0) \Delta \sigma_{ge} (\tau)}_{st} \,.
	\nonumber \\
\end{eqnarray}
%%
Mas adelante explicaremos el significado de esta funci\'on de dos tiempos. 
Cuando el intervalo de tiempo $\tau$ entre los dos operadores es cero 
se tiene que la intensidad de la emisi\'on es proporcional a
%
\begin{eqnarray} 
\vxp{\sigma_{eg} \sigma_{ge}}_{st} 
	&=& \vxp{\sigma_{ee}}_{st} = |\alpha_{eg} |^2   
	+\vxp{\Delta \sigma_{eg} \Delta \sigma_{ge}}_{st}  \,. 
\end{eqnarray}
%%
Decimos que el primer t\'ermino del lado derecho, $|\alpha_{eg} |^2$, 
es la parte coherente, dependiente \'unicamente del valor esperado del 
dipolo, y el segundo t\'ermino es la parte incoherente, 
%
\begin{eqnarray} 
\vxp{\Delta \sigma_{eg} \Delta \sigma_{ge}}_{st} 
	&=& \alpha_{ee} -|\alpha_{eg}|^2 		 \\  
	&=& \frac{ 2\Omega_0^4}
   	 { [ 2\Omega_0^2 +\gamma^2 +4\Delta^2 ]^2} \,.   \nonumber 
\end{eqnarray}
%%
El t\'ermino ``incoherente'' se refiere a que depende de las 
\textit{fluctuaciones}  de los operadores at\'omicos.

Una relaci\'on importante es el cociente entre la parte coherente, llamada 
tambi\'en el\'astica, y la incoherente, o inel\'astica: 
%
\begin{eqnarray}  
\frac{I_{coh}}{I_{inc}} &=& \frac{|\vxp{\sigma_{eg}}_{st}|^2}
	{\vxp{\Delta \sigma_{eg} \Delta \sigma_{ge}}_{st}}		\\ 
&=& \frac{ \Omega_0^2 [\gamma_a^2 +4\Delta^2]/D^2 }
	{ 2 \Omega_0^4 /D^2} 
	= \frac{ \gamma_a^2 +4\Delta^2 }{ 2 \Omega_0^2} \,, \nonumber \\ 
D &=& 2\Omega_0^2 +\gamma_a^2 +4\Delta^2 \,.
\end{eqnarray}
%%
Este es justamente el inverso del par\'ametro de saturaci\'on $I_S$ visto 
anteriormente, que es la intensidad de campo necesaria para saturar al 
\'atomo. Un valor $I_S <1$ indica que la intensidad incoherente es mayor 
que la coherente. 

La raz\'on de emisi\'on de fotones es $\gamma_a \vxp{\sigma_{ee}}$, que 
en el r\'egimen de saturaci\'on en un per\'iodo brillante es $\gamma_a/2$. 
El n\'umero de fotones emitidos en el intervalo $[0,t]$ es 
%
\begin{eqnarray} 
N(t) = \gamma_a \int_0^t \vxp{\sigma_{ee}(t')} \, dt' \,.
\end{eqnarray}
%%
Equivalentemente, podemos calcular el tiempo promedio entre la 
emisi\'on de dos fotones consecutivos, 
$\bar{\tau} = [\gamma_a \vxp{\sigma_{ee}}_{st}]^{-1}$. 
\textcolor{blue}{Comentar: funci\'on de retardo, etc.}


%%%%%%%%10%%%%%%%%20%%%%%%%%30%%%%%%%%40%%%%%%%%50%%%%%%%%60%%%%%%%%70%%%
\section{Probabilidad de Detecci\'on de Uno y Dos Fotones: 
Funciones de Correlaci\'on} 
Apunte en proceso
\subsection{Probabilidad de Detecci\'on de un Fot\'on} %%%%%%%%
Consideremos un detector en equilibrio t\'ermico a temperatura ambiente 
$T \sim 300$ K. La probabilidad de que los \'atomos sean excitados por la 
energ\'ia t\'ermica es muy peque\~na [de la distribuci\'on de Boltzmann 
$P_B \propto \exp{(-\hbar \omega/k_B T)} \sim \exp{(-40)} \to 0$]. Por lo 
tanto, los \'atomos inician en su estado inferior y solo absorben radiaci\'on, 
es decir, solo el operador de destrucci\'on juega un papel: 
%
\begin{eqnarray} 
 \hat{E}^+ (x) \ket{i} \propto a \ket{i} = \ket{f} \,.
\end{eqnarray}
%%

La probabilidad de transici\'on por absorci\'on de un fot\'on en la posici\'on 
$\mbr$ al tiempo $t$, en $x= \mbr, t$, es proporcional a 
%
\begin{eqnarray} 
W = W_{i \to f} &=& |\bra{f} \hat{E}^+ (x) \ket{i} |^2 	\nonumber \\
    &=& \bra{i} \hat{E}^- (x) \ket{f} \bra{f} \hat{E}^+ (x) \ket{i} ,
 \end{eqnarray}
%%
donde $\ket{i}$ y $\ket{f}$ son el estado inicial y final, respectivamente, 
del sistema acoplado \'atomo-campo. No nos interesa el estado final, de 
hecho, puede haber muchos de ellos, por lo que escribimos  
%
\begin{eqnarray}
W &=& \sum_f |\bra{f} \hat{E}^+ (x) \ket{i} |^2 
  = \sum_f \bra{i} \hat{E}^- (x) \ket{f} \bra{f} \hat{E}^+ (x) \ket{i} \nonumber \\ 
&=& \bra{i} \hat{E}^- (x) \left[ \sum_f \ket{f} \bra{f} \right] 
    \hat{E}^+ (x) \ket{i} 	\nonumber \\ 
    &=& \bra{i} \hat{E}^- (x)  \hat{E}^+ (x) \ket{i} . 
\end{eqnarray}
%%

Aunque sabemos que el \'atomo inicia en el estado base, normalmente 
no conocemos el estado inicial del campo, incluyendo variaciones 
estad\'isticas sobre $\ket{i}$. Usando el operador de densidad 
$\rho = \sum_i P_i \ketbra{i}{i}$ tenemos 
$W= \tr{\rho \hat{E}^- (x)  \hat{E}^+ (x)}$, que es la funci\'on de 
correlaci\'on de primer orden!

\subsection{Probabilidad de Detecci\'on de Dos Fotones} %%%%%%%%
Experimento de Hanbury-Brown y Twiss
%\newline Funciones de correlaci\'on 
%\newline F\'ormula cu\'antica de regresi\'on

La funci\'on de correlaci\'on $G^{(2)}$ se define operacionalmente como la 
probabilidad de detecci\'on de un fot\'on a tiempo $t$ y otro al tiempo 
$t+\tau$, 
%
\begin{eqnarray*}  	%\label{ec:G2operacional}
|\Psi(t+\tau) \rangle &=& A_{ge} (t+\tau) A_{ge} (t) |\Psi (0) \rangle  	\\
	&=& e^{iH(t+\tau)} A_{ge} e^{-iH(t+\tau)} e^{iHt} A_{ge} e^{-iHt} 
		|\Psi (0) \rangle \\ 
&=& e^{iH(t+\tau)} A_{ge} e^{-iH(t+\tau)} e^{iHt} A_{ge} |\Psi (t) \rangle \\ 
&=& e^{iH(t+\tau)} A_{ge} e^{-iH(t+\tau)} e^{iHt} |\Psi(t)=g \rangle \\ 
&=& e^{iH(t+\tau)} A_{ge} e^{-iH\tau}  |\Psi(t) \rangle \\ 
&=& e^{iH(t+\tau)} A_{ge} |\Psi(t+\tau) \rangle \\
&=&	e^{iH(t+\tau)} |\Psi(t+\tau)=g \rangle 
\end{eqnarray*} 	
%%
%
\begin{eqnarray*}  	\label{ec:G2operacional}
G^{(2)} (t,t+\tau) &=& \left| A_{ge} (t+\tau) A_{ge} (t) |\Psi \rangle  \right|^2 	
	\nonumber \\ 
&=& \langle \Psi(0) | A_{eg} (t) A_{eg} (t+\tau)  \\ 
	&& A_{ge} (t+\tau) A_{ge} (t) |\Psi(0) \rangle \\
&=& \langle  A_{eg} (t) A_{eg} (t+\tau)  \\ 
	&& A_{ge} (t+\tau) A_{ge} (t) |\Psi(0) \rangle \langle \Psi(0) |  \rangle \\
&=& \langle  A_{eg} (t) A_{eg} (t+\tau)  A_{ge} (t+\tau) A_{ge} (t)\rho(0) \rangle \\
&=& \langle  A_{eg} (t) A_{eg} (t+\tau)  A_{ge} (t+\tau) A_{ge} (t) \rangle \,.
\end{eqnarray*} 	
%%

La detecci\'on de un fot\'on viene asociada al decaimiento del \'atomo por 
emisi\'on espont\'anea, por lo que el \'atomo pasa del estado excitado al 
estado base, descrito por el operador $A_{0j}$, perdiendo la memoria de la 
historia previa; estad\'isticamente, sin embargo, todas las historias son 
equivalentes si las condiciones iniciales se repiten. De este modo el sistema 
alcanza el estado estacionario, en el cual las funciones a dos tiempos se 
reducen a funciones del intervalo entre ellos, 
$G^{(2)} (t,t+\tau) \to G^{(2)} (\tau)$,  donde 
%
\begin{eqnarray}
G^{(2)} (\tau) 
	= \langle A_{eg} (0) A_{eg} (\tau)  A_{ge} (\tau) A_{ge} (0) \rangle \,.
\end{eqnarray} 	
%%
Cuando $\tau \to \infty$ los fotones son completamente 
independientes, por lo que 
%
\begin{eqnarray}
\langle A_{eg} (0) A_{eg} (\tau)  A_{ge} (\tau) A_{ge} (0) \rangle &\to&  
	\langle A_{eg} A_{ge} \rangle(0) \langle A_{eg} A_{ge} \rangle(\tau) 	\nonumber \\ 
	&=& \langle A_{ee} \rangle_{st}^2 \,, 
\end{eqnarray} 	
%%
lo cual nos sirve para normalizar la correlaci\'on,
%
\begin{eqnarray}
g^{(2)} (\tau) 
	&=& \frac{ \langle A_{eg} (0) A_{eg} (\tau)  A_{ge} (\tau) A_{ge} (0) \rangle }
	{ \langle A_{ee} \rangle_{st}^2 } \,, 
\end{eqnarray} 	
%%
tal que  
%
\begin{eqnarray} 
g^{(2)} (\tau \to \infty) = 1 \,. 
\end{eqnarray} 	
%%
%
\begin{eqnarray*}
G^{(2)} (\tau) 
	&=& \langle A_{eg} (0) A_{eg} (\tau)  A_{ge} (\tau) A_{ge} (0) \rangle 	\\
	&=&  \langle A_{eg} (\tau)  A_{ge} (\tau) \rangle_{\rho(0)=\rho_{gg}} \\ 
	&=& \langle A_{ee} (\tau) \rangle_{\rho(0)=\rho_{gg}}
\end{eqnarray*} 	
%%

Usando la f\'ormula cu\'antica de regresi\'on la correlaci\'on $g_j^{(2)}$ se 
puede reducir a la evoluci\'on de la poblaci\'on del estado excitado dado que 
inici\'o en el estado base, normalizada por la poblaci\'on promedio 
\cite{Carm99}, 
%
\begin{eqnarray}
g^{(2)} (\tau) &=& \frac{ \langle A_{ee} (\tau)\rangle_{\rho(0)=\rho_{gg}} }
	{ \langle A_{ee} \rangle_{st} }
\end{eqnarray}
%%
es decir, la correlaci\'on se reduce a normalizar la poblaci\'on del estado 
excitado. Esto \'ultimo, sin embargo, no ocurre para todos los sistemas 
\'opticos (p. ej., fotones que emergen de una cavidad con \'atomos) \ldots 


%%%%%%%%10%%%%%%%%20%%%%%%%%30%%%%%%%%40%%%%%%%%50%%%%%%%%60%%%%%%%%70%%%
\section{Excitaci\'on Con Campos Dependientes del Tiempo} 
teorema del \'area; modelo LZ, STIRAP ?, etc

%%5%%%%%%%%60%%%%%%%%70%%%
\subsection{Teorema del \'Area}
La f\'ormula de Rabi, 
%
\begin{eqnarray} 	\label{eq:RabiFormula}
|C_a (t)|^2 &=& \frac{\Omega_0^2}{\Omega^2} \sin^2{(\Omega t/2)}  \,,  
\end{eqnarray}
%% 
que expresa la probabilidad de transici\'on en un \'atomo de dos niveles, 
no se limita a excitaci\'on monocrom\'atica. Un caso de especial inter\'es 
ocurre cuando cuando los l\'aseres son pulsados. En este caso decimos 
que el pulso tiene un \textit{\'area}  
$A = \int_{-\infty}^{\infty} \Omega_0 (t') dt' $. Si $A=\pi$ el \'atomo 
experimenta una inversi\'on de la probabilidad: un \'atomo en estado 
base es promovido a su estado excitado, y viceversa. Un pulso de \'area 
$A=\pi/2$ crea una superposici\'on de amplitudes iguales entre los estados. 
Este m\'etodo es usado rutinariamente en muchas aplicaciones, m\'as 
notablemente en procesamiento cu\'antico de informaci\'on.   

Efecto de la forma de los pulsos \ldots 

Secuencias de pulsos de diferentes \'areas: eco de fotones,  

%%5%%%%%%%%60%%%%%%%%70%%%
\subsection{Seguimiento Adiab\'atico}
En algunas situaciones el m\'etodo de inversi\'on de poblaci\'on sugerido 
por el teorema del \'area no es robusto: requiere de alta uniformidad en la 
amplitud del pulso y exactitud en la resonancia $\Delta =0$ del centro del 
espectro del pulso con la frecuencia de transici\'on. Un m\'etodo mucho 
mas efectivo para crear inversi\'on de poblaci\'on utiliza campos con 
frecuencia variable (adiabatic following), ya sea con amplitud constante 
o pulsado.  

\textcolor{red}{Ilustrar con los estados vestidos, Secc.~\ref{dressedST} !!!}

Un modelo particular es el de Landau y Zener $\Omega_0 =cte$, 
$\Delta = \beta^2 t$, donde $\beta^2$ es la pendiente de la desinton\'ia en 
resonancia. Landau y Zener, independientemente, resolvieron el problema. 
Si $-\infty < t < \infty$ encontraron que 
$P_e (t \to \infty) = 1- \exp{ (-\pi \Omega_0^2/2\beta^2) }$. La duraci\'on 
del salto alrededor de la resonancia hacia el primer m\'aximo es 
$\tau_0 = \omega/\beta^2$. Si $\Omega > 3\beta$ la inversi\'on es total, 
$P_e (t \to \infty) \to 1$. Por supuesto, hemos ignorado de momento que 
hay decaimientos, por emisi\'on espont\'anea, por ejemplo. Sin embargo, 
los resultados son robustos si el proceso ocurre en un tiempo mucho 
menor que el de decaimiento, lo cual es factible; en realidad, el proceso 
no requiere un tiempo muy largo. La soluci\'on de Landau y Zener es 
idealizada pero muy efectiva y compacta. Existen otros modelos de 
desinton\'ia y frecuencia de Rabi con dependientes del tiempo lenta dan 
resultados similares.  


%%5%%%%%%%%60%%%%%%%%70%%%
\subsection{STIRAP}
Ver Unidad sobre efectos en \'atomos multinivel. 


%%%%%%%%10%%%%%%%%20%%%%%%%%30%%%%%%%%40%%%%%%%%50%%%%%%%%60%%%%%%%%70%%%
\section{Coherencia Cu\'antica} 
Sea $G^{(m)}(x_1 \ldots x_m; y_m \ldots y_1)$ la funci\'on de coherencia 
cl\'asica de orden $m$, 
%
\begin{eqnarray*}
\lefteqn{G^{(m)}(x_1 \ldots x_m; y_m \ldots y_1) } \hspace{1cm}	\\
&=& \langle E^-(x_1) \cdots E^-(x_m) E^+(y_m) \cdots E^+(y_1) \rangle \,, 
\end{eqnarray*}
%% 
donde $x_i,y_i =\mathbf{r}_i t_i$. Un campo es coherente a orden $m$ si 
$G^{(m)}$ es factorizable 
%
\begin{eqnarray*}
\lefteqn{ G^{(m)}(x_1 \ldots x_m; y_m \ldots y_1) } \hspace{1cm}	\\
&=&  \mathcal{E}^{\ast}(x_1) \cdots \mathcal{E}^{\ast}(x_m) 
	\mathcal{E} (y_m) \cdots \mathcal{E} (y_1) \,. 
\end{eqnarray*}
%% 
En la teor\'ia cu\'antica de la coherencia reemplazamos $E^{\pm}(x_i)$ 
por operadores $\hat{E}^{\pm}(x_i)$, y los promedios 
$\langle \cdots \rangle$ son ahora promedios cu\'anticos (valores 
esperados): 
%
\begin{eqnarray*}
\lefteqn{ G^{(m)}(x_1 \ldots x_m; y_m \ldots y_1) } \hspace{1cm}	\\
&=& \langle \hat{E}^-(x_1) \cdots \hat{E}^-(x_m) \hat{E}^+(y_m) 
    \cdots \hat{E}^+(y_1) \rangle  \\ 
&=& \tr{ \rho \hat{E}^-(x_1) \cdots \hat{E}^-(x_m) \hat{E}^+(y_m) 
    \cdots \hat{E}^+(y_1) } . 
\end{eqnarray*}
%% 

Como ejemplos de campo EM consideremos una onda viajera unimodal en 
espacio libre de amplitud 
$\mcE_{\omega} = \sqrt{\hbar \omega /\epsilon_0 V}$,    
$$ \hat{E}(\mbr,t) = \mcE_{\omega} a e^{-i\omega t} f(\mbr) 
+\mcE_{\omega} \adg e^{i\omega t} f^\ast (\mbr) , $$
donde $\omega$ es la frecuencia angular, V es el volumen y $f(\mbr)$ 
describe la estructura espacial del modo. Para una onda estacionaria en 
una cavidad unidimensional 
%
\begin{eqnarray*}
\hat{E}^+(x) &=& \mcE_{\omega} a e^{-i\omega t} \sin{kz} ,    \\ 
\hat{E}^-(x) &=& \left[ \hat{E}^+(x) \right]^+  
   = \mcE_{\omega} \adg e^{i\omega t} \sin{kz} , 
\end{eqnarray*}
%% 
entonces 
%
\begin{eqnarray*}
G^{(1)}(x_1,x_1) &=& \mcE_{\omega}^2 \sin^2{kz} \,\, \tr{(\rho \adg a)} \\ 
G^{(1)}(x_1,x_2) &=& \mcE_{\omega}^2 \sin^2{kz_1} \sin^2{kz_2} 
    e^{i\omega (t_1-t_2)} \tr{(\rho \adg a)}
\end{eqnarray*}
%% 
Puesto que $\tr{(\rho \adg a)} = \bar{n}$ la correlaci\'on de primer 
orden es \text{insensible} a la estad\'istica! Se requieren correlaciones 
de orden superior para distinguir fuentes con diferentes estad\'isticas. 

Consideremos ahora efectos del estado del campo EM en las correlaciones: 
1) Campo en estado coherente 
($a \ket{\alpha} = \alpha \ket{\alpha}$, 
$\bra{\alpha} \adg = \bra{\alpha} \alpha^\ast$ ) 
%
\begin{eqnarray*}
\lefteqn{ G^{(m)}(x_1 \ldots x_m; y_m \ldots y_1) } \hspace{2cm} \\
&\propto& 
\bra{\alpha} \adg (x_1) \cdots \adg (x_m) a(y_m) \cdots a(y_1) 
    \ket{\alpha}  \\ 
&=& \mcE^\ast (x_1) \ldots \mcE^\ast (x_m) \mcE (y_m) \ldots \mcE (y_1) , 
\end{eqnarray*}
%% 
donde $\mcE (x_i) = \alpha e^{-i \omega t}$. La correlaci\'on de segundo orden normalizada es  
%
\begin{eqnarray*}
g^{(2)}(x_1,x_2; x_2,x_1) &=& \frac{ G^{(2)} (x_1,x_2; x_2,x_1) }
    {G^{(1)} (x_1) G^{(1)} (x_1)} \\ 
  &=& \frac{ \mcE^\ast(x_1) \mcE^\ast(x_2) \mcE (x_2) \mcE (x_1) }
    { \mcE^\ast(x_1) \mcE (x_1) \mcE^\ast(x_2) \mcE (x_2) } = 1 , 
\end{eqnarray*}
%% 
esto es, el estado coherente de Glauber es coherente a todos los 
\'ordenes!

2) Campo en estado de Fock (n\'umero) $\ket{n}$, $\rho =\ketbra{n}{n}$ 
(estado puro)
%
\begin{eqnarray*}
G^{(1)}(x_1) &=& \bra{n} \hat{E}^- (x_1) \hat{E}^+ (x_1) \ket{n}   \\ 
&=& |\mcE_{\omega} \sin{kz}|^2 \bra{n} \adg a \ket{n} 
    = n |\mcE_{\omega} \sin{kz}|^2 ,     \\ 
|G^{(1)}(x_1)|^2 &=& n^2 |\mcE_{\omega} \sin{kz}|^4 . 
\end{eqnarray*}
%% 
%
\begin{eqnarray*}
G^{(2)}(x_1,x_2) 
&=& \bra{n} \hat{E}^- (x_1) \hat{E}^- (x_1) \hat{E}^+ (x_1) 
    \hat{E}^+ (x_1) \ket{n}   \\ 
&=& |\mcE_{\omega} \sin{kz}|^4 \bra{n} \adg \adg a a \ket{n} \\  
&=& n(n-1) |\mcE_{\omega} \sin{kz}|^4 = G^{(2)} (0) .
   \end{eqnarray*}
%% 
Puesto que $G^{(2)}(x_1,x_2) \neq |G^{(1)}(x_1)|^2$ la correlaci\'on no 
se factoriza. Entonces un estado de Fock no es coherente a segundo orden 
(ni a \'ordenes superiores). Veamos la normalizaci\'on a $x=0$, 
%
\begin{eqnarray*}
g^{(2)}(0) &=& \frac{ G^{(2)} (0) }{ |G^{(1)} (0)|^2 } \\ 
  &=& \frac{n(n-1)}{n} = 1 - \frac{1}{n} . 
\end{eqnarray*}
%% 
El estado de Fock presenta antibunching, $g^{(2)} < 1$, pero se aproxima 
al l\'imite cl\'asico $g^{(2)} \to 1$ para $n \to \infty$. 

\subsection{Campo EM en Equilibrio T\'ermico} %%%%%%%%%%%%%%%
(P'aginas 9 y 10 de viejo apunte manuscrito)
\newline Conclu\'imos la clase pasada escribiendo la forma general de 
un campo EM como una superposici\'on de estados de una base completa 
$\ket{n_q}$ ($q= \mbk \lambda$), 
$$ \ket{\Psi} = \sum_q C_{\{n_q\}} \ket{\{n_q\}} .$$ 
Los estados coherentes y comprimidos pueden escribirse de esta forma,  
Sin embargo, el campo t\'ermico (CT) no puede escribirse en t\'erminos 
de estados puros. como $\ket{\Psi}$ arriba. Un CT es una superposici\'on 
\textit{incoherente} de ondas con amplitud y fase definidas, por lo que cada 
onda contribuye en t\'erminos probabil\'isticos. Cuando \'unicamente 
se puede especificar un conjunto de probabilidades de que el campo se 
encuentre en un rango de estados, cada uno corresponde a a un estado de 
un conjunto completo, entonces el estado se encuentra en una 
\textit{mezcla estad\'istica} (no una superposici\'on). En QM estas 
mezclas est\'an dadas por el operador de densidad. 

Consideremos un campo en una cavidad en el que conocemos la probabilidad 
$P_R$ ($\sum_R P_R$) de que el campo est\'e en el estado $\ket{R}$ (p. ej., 
en un estado $\ket{n}$),   
%
\begin{eqnarray*}
\vxp{O} &=& \bra{R} O \ket{R}  \qquad \qquad \text{estados puros} \\ 
\vxp{O} &=& \sum_R P_R \bra{R} O \ket{R} \qquad \text{mezcla estad.} 
\end{eqnarray*}
%%

Sea $\ket{s}$ una base completa de estados puros. Por el teorema de la 
cerradura, $\sum_s \ketbra{s}{s} =1$, 
%
\begin{eqnarray*}
\vxp{O} &=& \sum_R P_R \bra{R} O \ketbra{s}{s} \ket{R}  
    = \sum_R P_R \braket{s}{R} \bra{R} O \ket{s}  \\  
    &=& \sum_s  \bra{s} \rho O \ket{s} = \tr{\rho O} , 
\end{eqnarray*}
%%
donde $\rho = \sum_R P_R \ketbra{R}{R}$ es el operador de densidad 
que contiene exactamente la misma informaci\'on que $P_R$, y $\rho$ 
esta determinada una vez que $P_R$ est\'e determinada para la base de 
estados $\ket{R}$. La traza es la suma de los elementos diagonales de 
la matriz para cualquier conjunto completo de estados. Esta 
expresi\'on es mejor que $\vxp{O} = \sum_R P_R \bra{R} O \ket{R}$. 
Si $O= \mathbf{1}$ entonces $\tr{\rho} =1, $
%
\begin{eqnarray*}
\bra{R'} \rho \ket{R''} 
    &=& \sum_R P_R \braket{R'}{R} \braket{R}{R''}     \\ 
    &=& P_{R'} \delta_{R' \,R''} 
\end{eqnarray*}
%%
solo los elementos diagonales son diferentes de cero, 
%
\begin{eqnarray*}
\bra{s} \rho \ket{s'} = \sum_R P_R \braket{s}{R} \braket{R}{s'} . 
\end{eqnarray*}
%%
Para un estado puro $\rho = \ketbra{R}{R} \to \rho^2 
= \ket{R} \braket{R}{R} \bra{R}  = \rho$. Nota: 
$\rho = \sum_n \sum_m \ketbra{n}{n} \rho \ketbra{m}{m} 
    = \sum_{n,m} \rho_{nm} \ketbra{n}{m}$, 

Un importante ejemplo de mezcla estad\'istica esta dado por la 
excitaci\'on t\'ermica de fotones de un modo de una cavidad a temperatura 
$T$, $P_n = \bar{n}^n /[(1+\bar{n})^{1+n}]$. Entonces 
%
\begin{eqnarray*}
\rho &=& \sum_n P_n \ketbra{n}{n} 
    = \sum_n \frac{\bar{n}^n}{(1+\bar{n})^{1+n}} \ketbra{n}{n} .
\end{eqnarray*}
%%

Calculemos los valores esperados del campo y del n\'umero de 
fotones de un CT: 
%
\begin{eqnarray*}
\vxp{a} &=& \tr{a \rho} = P_n \sum_n \bra{n} a \ket{n}  \\ 
&=& \sum_n P_n \sqrt{n} \braket{n}{n-1} = 0 . 
\end{eqnarray*}
%%
%
\begin{eqnarray*}
\vxp{\adg a} &=& \tr{\adg a \rho} 
    = P_n \sum_n \bra{n} \adg a \ket{n}  \\ 
&=& \sum_n n P_n \braket{n}{n} = \bar{n} . 
\end{eqnarray*}
%%
Para un CT multimodal 
%
\begin{eqnarray*}
\vxp{\adg_q a_q} &=& \Pi_q \frac{\bar{n}^n_q}{(1+\bar{n})^{1+n_q}} \\ 
\rho &=& \sum_q P_{n_q} \ketbra{\{n_q\}}{\{n_q\}} .
\end{eqnarray*}
%%


%%%%%%%%10%%%%%%%%20%%%%%%%%30%%%%%%%%40%%%%%%%%50%%%%%%%%60%%%%%%%%70%%%
\section{Soluci\'on Formal Matricial de las Ecuaciones de Bloch y 
Correlaciones a Dos Tiempos} 
En esta secci\'on mostraremos el m\'etodo matricial para resolver las 
ecuaciones de Bloch y correlaciones a dos tiempos, en preparaci\'on 
para soluciones num\'ericas. Los detalles de las soluciones se encuentran 
en notas aparte. Tomamos el caso particular del \'atomo de dos niveles 
interactuando con un l\'aser monocrom\'atico.  

Escribimos las ecuaciones de Bloch en forma matricial,  
%
\begin{eqnarray} 	\label{eq:SBloch}
\langle \dot{\mathbf{s}} \rangle = 
	\mathbf{M} \langle \mathbf{s} \rangle +\mathbf{b}  \,,	\\ 
%
\frac{d}{dt} \left( \langle \mathbf{s} \rangle 
	+\mathbf{M}^{-1} \mathbf{b} \right) 
&=& \mathbf{M} \left( \langle \mathbf{s} \rangle 
	+\mathbf{M}^{-1}\mathbf{b} \right) 	\,, 	\nonumber 
\end{eqnarray}
%%
donde $\mathbf{s} = (\sigma_-, \sigma_+,\sigma_z)^T$ es un vector que 
contiene los operadores del pseudoesp\'in 
$\sigma_{jk} = |j \rangle \langle k|$ que representan las poblaciones 
($j=k$) y las coherencias ($j \neq k$), $\mathbf{M}$ es la matriz de los 
coeficientes, y $\mathbf{b}$ es un vector diferente de cero cuando el 
n\'umero de ecuaciones es menor que el de poblaciones m\'as coherencias. 

La soluci\'on formal puede escribirse como 
%
\begin{subequations}
\begin{eqnarray}
\langle \mathbf{s}(t) \rangle &=& e^{ \mathbf{M} t } \left[ 
	 \langle \mathbf{s}(0) \rangle +\mathbf{M}^{-1} \mathbf{b} \right] 
	-\mathbf{M}^{-1} \mathbf{b}	\\
&=& S^{-1} e^{\Lambda t} S \left[ \langle \mathbf{s}(0) \rangle 
	+\mathbf{M}^{-1} \mathbf{b} \right] -\mathbf{M}^{-1} \mathbf{b} \,, 
\end{eqnarray} 
\end{subequations}
%%
donde $S$ es una matriz que diagonaliza $\mathbf{M}$, y $\Lambda$ 
es la matriz diagonal con los valores propios $\lambda_m$.  

La soluci\'on en estado estacionario ($t \to \infty$) de las ecuaciones de 
Bloch es 
%
\begin{equation} 	\label{eq:Bloch_est}
0 =\mathbf{M} \langle \mathbf{s} \rangle_{st} +\mathbf{b}  \,, 
    \qquad \to \langle \mathbf{s} \rangle_{st} 
       =\mathbf{M}^{-1} \mathbf{b} \,.  	
\end{equation}
%%

%%%%%%%%10%%%%%%%%20%%%%%%%%30%%%%%%%%40
\subsection{Correlaciones a Dos Tiempos} 
La f\'ormula cu\'antica de regresi\'on permite escribir ecuaciones para 
funciones de correlaci\'on a dos (o m\'as) tiempos de la misma forma 
que para un tiempo: 
%
\begin{eqnarray*}
\frac{d}{d \tau} \langle \sigma_{+}(0) \mathbf{s} (\tau) \rangle_{st} 
 &=& \mathbf{M} \langle \sigma_{+}(0) \mathbf{s} (\tau) \rangle_{st} 
	+\mathbf{b} \langle \sigma_{+} \rangle_{st} \,,	\\ 
%
\frac{d}{d \tau} \langle \mathbf{s} (\tau) \sigma_{-}(0)  \rangle_{st} 
 &=& \mathbf{M} \langle \mathbf{s} (\tau) \sigma_{-}(0)  \rangle_{st} 
	+\mathbf{b} \langle \sigma_{-} \rangle_{st} \,,	\\ 
%
\frac{d}{d \tau} \langle \sigma_{+}(0) \mathbf{s} (\tau) 
	\sigma_{-}(0) \rangle_{st}  
 &=& \mathbf{M} \langle \sigma_{+}(0) \mathbf{s} (\tau) 
 	\sigma_{-}(0) \rangle_{st}  	\\ 
	&& +\mathbf{b} \langle \sigma_{+} \sigma_{-} \rangle_{st} 	\,,
\end{eqnarray*}
%%

Veamos una forma compacta de estas ecuaciones, 
%
\begin{eqnarray}
\frac{d}{d\tau} \langle A_+(0) \mathbf{s}(\tau) A_-(0) \rangle_{st} 
	&=& \mathbf{M} \langle A_+(0) \mathbf{s}(\tau) A_-(0) \rangle_{st} 
	\nonumber \\ 
	&& +\langle A_+ A_- \rangle_{st} \mathbf{b}	\,.   	
\end{eqnarray}
%%
La FCR debe cumplir algunas condiciones pero nos limitaremos a un 
procedimiento sencillo: (i) remover los brackets a 
$\langle \mathbf{s} \rangle$, (ii) multiplicar $\mathbf{s}$ por la izquierda y 
por la derecha por los operadores $A_+$  and $A_-$, respectivamente, 
(iii) cambiar $t$ por $\tau$, y (iv) poner de nuevo los brackets. La soluci\'on 
es formalmente id\'entica a la de un tiempo, 
%
\begin{eqnarray}
\lefteqn{  \langle A_+(0) \mathbf{s}(\tau) A_-(0) \rangle_{st} 
	= -\langle A_+ A_- \rangle_{st} \mathbf{M}^{-1} \mathbf{b} } 
	\hspace{20mm}	 	\nonumber \\
&& +e^{ \mathbf{M} \tau } \left[ \langle A_+ \mathbf{s} A_- \rangle_{st} 
	 +\langle A_+ A_- \rangle_{st} \mathbf{M}^{-1} \mathbf{b} \right] \,. 
	 \nonumber \\
\end{eqnarray}
%%
Similarmente, podemos reemplazar $e^{ \mathbf{M} \tau }$ con 
$S^{-1} e^{\Lambda \tau} S$. 

\subsection{Correlaciones de Fotones} 	%%%%%%%%%%%

Una de las funciones de correlaci\'on m\'as importantes es la que describe 
la probabilidad de que el \'atomo emita un fot\'on a tiempo $\tau=0$ y otro 
a tiempo $\tau$,  
%
\begin{eqnarray}
G^{(2)}  (\tau) &=& \langle \sigma_+(0) \sigma_+ (\tau) \sigma_- (\tau)  
	\sigma_-(0) \rangle_{st}  	\\
 &=& \frac{1}{2} \left[ \langle \sigma_+ \sigma_- \rangle_{st} 
 	+ \langle \sigma_+(0) \sigma_+ \sigma_z (\tau)  \sigma_-(0) \rangle_{st}
 	\right] \,, 	\label{eq:G2_tau} \nonumber \\
 \end{eqnarray}
%%
donde usamos $\sigma_+ \sigma_- = (1+\sigma_z)/2$. Resolveremos 
entonces 
%
\begin{eqnarray}
\frac{d}{d \tau} \langle \sigma_{+}(0) \mathbf{s} (\tau) 
	\sigma_{-}(0) \rangle_{st}  	
 &=& \mathbf{M} \langle \sigma_{+}(0) \mathbf{s} (\tau) 
 	\sigma_{-}(0) \rangle_{st}  		\nonumber 	\\ 
	&& +\mathbf{b} \langle \sigma_{+} \sigma_{-} \rangle_{st} 	\,,  
\end{eqnarray}
%%
cuya soluci\'on formal es 
%
\begin{eqnarray}
\lefteqn{  \langle \sigma_{+}(0) \mathbf{s} (\tau) \sigma_{-}(0) \rangle_{st}  
	= -\langle \sigma_+ \sigma_- \rangle_{st} \mathbf{M}^{-1} \mathbf{b} } 
	\hspace{20mm}	 	\nonumber \\
&& +e^{ \mathbf{M} \tau } \left[ \langle \sigma_+ \mathbf{s} \sigma_- \rangle_{st} 
	 +\langle \sigma_+ \sigma_- \rangle_{st} \mathbf{M}^{-1} \mathbf{b} \right] \,, 
	 \nonumber \\
\end{eqnarray}
%%
y sus condiciones iniciales son
%
\begin{equation}
\langle \sigma_+ \mathbf{s} \sigma_- \rangle_{st} 
	= \langle \sigma_+ \sigma_- \rangle_{st} (0,0,-1)^T \,.
\end{equation}
%% 
Ahora podemos escribir la soluci\'on como
%
\begin{eqnarray}
\lefteqn{  \langle \sigma_{+}(0) \mathbf{s} (\tau) \sigma_{-}(0) \rangle_{st}  
	= \langle \sigma_+ \sigma_- \rangle_{st} 
	\left\{ - \mathbf{M}^{-1} \mathbf{b} \right. 	} 
	\hspace{20mm}	 	\nonumber \\
&& \left. +e^{ \mathbf{M} \tau } \left[ (0,0,-1) 
	 + \mathbf{M}^{-1} \mathbf{b} \right] 	\right\} \,, 	\nonumber \\ 
%
&=& \langle \sigma_+ \sigma_- \rangle_{st} 
	\langle \mathbf{s} (\tau) \rangle_{\rho(0)= |g \rangle \langle g| } \,.
\end{eqnarray}
%%
Notemos que $(0,0,-1)^T$ es simplemente la condici\'on inicial 
$\mathbf{s} (0)$ para un \'atomo preparado en su estado base, 
$\rho(0) = |g \rangle \langle g|$. N\'otese que en este c\'alculo es 
independiente de si $\Delta$ es cero o no en $\mathbf{M}$, a diferencia 
del caso que veremos sobre el espectro.Sustituyendo el tercer elemento 
de la ecuaci\'on anterior [$\mathbf{s}|_3 = \sigma_z$] en la 
Ec.~(\ref{eq:G2_tau}) establece nuestro resultado: 
%
\begin{eqnarray}
G^{(2)}  (\tau) &=& \langle \sigma_+ \sigma_- \rangle_{st} \frac{1}{2} 
	\left[ 1+ \langle \sigma_z (\tau) \rangle_{\rho(0)= |g \rangle \langle g| }
 	\right] \,,  \nonumber \\ 
&=& \rho_{ee}^{st} \rho_{ee}(\tau)_{\rho(0)= |g \rangle \langle g| } \,, \\ 
%
g^{(2)}  (\tau) &=& \frac{ \rho_{ee}(\tau)_{\rho(0)= |g \rangle \langle g| }}
	{\rho_{ee}^{st}} \,.
 \end{eqnarray}
%%
Mencionamos anteriormente que no es posible una soluci\'on anal\'itica 
general a todo tiempo; para el caso de $\Delta=0$ la correlaci\'on de dos 
fotones es   
%
\begin{eqnarray}
g^{(2)}  (\tau) &=& 1- e^{-3\gamma \tau /4} \left( \cosh{\delta \tau} 
	-\frac{3\gamma/4}{\delta} \sinh{\delta \tau} \right) \,, \nonumber \\
 \end{eqnarray}
%% 
donde $\delta = \sqrt{(\gamma/4)^2 -\Omega_0^2}$. 

%%%%%%%%10%%%%%%%%20%%%%%%%%30%%%%%%%%40
\subsection{Operadores de Fluctuaciones}
Para las ecuaciones de Bloch, $\langle \mathbf{\dot{s}} \rangle = 
\mathbf{M} \langle \mathbf{s} \rangle +\mathbf{b}$, vimos que el estado 
estacionario ($t \to \infty$) es 
$\langle \mathbf{s} \rangle_{st} =\mathbf{M}^{-1} \mathbf{b}$. Dividimos 
el campo del dipolo en su media m\'as fluctuaciones, 
$\mathbf{s}(t) = \langle \mathbf{s} \rangle_{st} +\Delta \mathbf{s}(t)$  
lo cual, sustituyendo en las ecuaciones de Bloch, nos lleva a  
%
\begin{eqnarray}
\frac{d}{dt} \langle \Delta \mathbf{s} \rangle +0
&=& \mathbf{M} \left[ \langle \mathbf{s} \rangle_{st} 
	+\langle \Delta \mathbf{s} \rangle \right] +\mathbf{b}  \nonumber \\
&=& \mathbf{M} \langle \Delta \mathbf{s} \rangle 
	+\mathbf{M} \langle \mathbf{s} \rangle_{st} +\mathbf{b}  
	\nonumber 	\\
\frac{d}{dt} \langle \Delta \mathbf{s} \rangle 
	&=& \mathbf{M} \langle \Delta \mathbf{s} \rangle  \,.
\end{eqnarray}
%%
Por supuesto, esta expresi\'on no tiene sentido porque, por definici\'on, 
$\langle \Delta \mathbf{s} \rangle =0$. Debemos recurrir a la FCR para 
calcular correlaciones a dos tiempos: (i) remover los brackets, multiplicar  
$\Delta \mathbf{s}$ por la izquierda y por la derecha por operadores 
apropiados $\Delta A_+$ y $\Delta A_-$, respectivamente, y (iii) colocar 
de nuevo los brackets, obteniendo 
%
\begin{eqnarray}
\frac{d}{d\tau} \langle \Delta A_+ \Delta \mathbf{s}(\tau) \Delta A_- \rangle 
 &=& \mathbf{M} \langle \Delta A_+ \Delta \mathbf{s}(\tau) \Delta A_- \rangle  \,, 
	\nonumber \\
\end{eqnarray}
%%
con soluciones 
%
\begin{eqnarray*}
\langle \Delta A_+ \Delta \mathbf{s}(\tau) \Delta A_- \rangle 
 &=& e^{\mathbf{M} \tau} \langle \Delta A_+ \Delta \mathbf{s} 
 	\Delta A_- \rangle  	\\ 
&=& S^{-1} e^{\Lambda t} S  \langle \Delta A_+ \Delta \mathbf{s} 
 	\Delta A_- \rangle 	\,.	
\end{eqnarray*}
%%
Notemos que las ecuaciones para las correlaciones de operadores de 
fluctuaci\'on forman un sistema homog\'eneo \ldots

%%%%%%%%10%%%%%%%%20%%%%%%%%30%%%%%%%%40
\subsection{Espectro Incoherente} 
Para la parte incoherente del espectro, 
%
\begin{eqnarray} 	%\label{eq:}
S_{inc}(\omega) 
&=& \frac{1}{ \pi \alpha_{ee} } \mathrm{Re} \left[  \int_0^{\infty} 
	d\tau e^{-i \omega \tau} 	
	\langle \Delta \sigma_+ (0 )\Delta \sigma_- (\tau) \rangle_{st} \right] \,,  
	\nonumber \\ 
\end{eqnarray}
%%
requirimos resolver la correlaci\'on para 
$\langle \Delta \sigma_+(0) \Delta \sigma_-(\tau) \rangle_{st}$. Sin 
embargo, por lo general, integramos de manera formal usando matrices
%
\begin{eqnarray} 	%\label{eq:}
S_{inc}(\omega) 
&=& \frac{1}{ \pi \alpha_{ee}  } \mathrm{Re} \left[  \int_0^{\infty} 
	d\tau e^{-(i \omega \mathbf{1} -\mathbf{M}) \tau} 	
	\langle \Delta \sigma_+ \Delta \mathbf{s} \rangle_{st} \right]_1   
	\nonumber \\ 
&=& \frac{1}{ \pi \alpha_{ee}  } \mathrm{Re} \left[ 
	-(i \omega \mathbf{1} -\mathbf{M})^{-1} 
	e^{-(i \omega \mathbf{1} -\mathbf{M}) \tau} \left. \right|_0^{\infty}  
	 \right. 	\nonumber \\
&&	\left. \times 
	\langle \Delta \sigma_+ \Delta \mathbf{s} \rangle_{st} \right]_1 	
	\nonumber \\ 
&=& \frac{1}{ \pi \alpha_{ee}  } \mathrm{Re} \left[ 
	(i \omega \mathbf{1} -\mathbf{M})^{-1} 
	\langle \Delta \sigma_+ \Delta \mathbf{s} \rangle_{st} \right]_1 	\,, 	
	\nonumber \\ 
\end{eqnarray}
%%
donde $\mathbf{1}$ es la matriz identidad. El sub\'indice $1$ fuera de los 
corchetes indica que solamente tomamos el primer elemento del vector 
resultante. 

Para apreciar mejor este m\'etodo notemos que la soluci\'on de la 
correlaci\'on $\langle \Delta \sigma_+ (0 )\Delta \sigma_- (\tau) \rangle_{st}$ 
se puede escribir como una suma de exponenciales, 
%
\begin{eqnarray}
\langle \Delta \sigma_+(0) \Delta \sigma_-(\tau) \rangle_{st} 
	= \sum_m C_m e^{\lambda_m \tau} \,, 
\end{eqnarray}
%%
donde $\lambda_m$ son los valores propios de $\mathbf{M}$, con parte 
real negativa. Entonces tenemos integrales tales como 
%
\begin{eqnarray} 	%\label{eq:integrals}
\int_0^{\infty} d\tau e^{(-i \omega +\lambda_m) \tau} 
	&=& \left. \frac{e^{(-i \omega +\lambda_m) \tau}}
	{-i \omega +\lambda_m} \right|_0^{\infty} 
	= \frac{-1}{-i \omega +\lambda_m} 	\,, \nonumber 	\\ 
\mathrm{Re} \int_0^{\infty} d\tau e^{(-i \omega +\lambda_m) \tau} 
	&=& - \mathrm{Re} \left[ \frac{i \omega +\lambda_m^\ast}
	{ \omega^2 +|\lambda_m|^2} \right] 
	= \frac{- \lambda_m}{ \omega^2 +\lambda_m^2} \,, \nonumber \\
\end{eqnarray}
%%
donde los picos tienen forma Lorentziana. \textcolor{red}{REVISAR el caso 
cuando $\lambda_m$ es complejo. ?`Qu\'e esta mal?}

%%%%%%%%10%%%%%%%%20%%%%%%%%30%%%%%%%%40
\subsection{Manipulando las Integrales - Espectro Integrado} 
El \'area debajo de la Lorentziana es una medida de la intensidad de un 
pico espectral. Al integrar sobre todas las frecuencias obtenemos 
%
\begin{eqnarray} 
\int_{-\infty}^{\infty} \frac{d\omega}{\omega^2 +a^2} 
	&=& \left. \frac{1}{a} \tan^{-1}{\left( \frac{\omega}{a} \right)} 
		\right|_{-\infty}^{\infty} 
	=  \frac{1}{a} \left[ \frac{\pi}{2} -\frac{-\pi}{2} \right] 	\nonumber \\ 
	&=& \pi/a	\,.
\end{eqnarray}
%%

%%%%%%%%10%%%%%%%%20    REFERENCES    40%%%%%%%%50%%%%%%%%60%%%%%%%%70%%%
\begin{references} 
%
\bibitem{Carm93} H. J. Carmichael, \textit{An Open Systems Approach to 
	Quantum Optics}, Lecture Notes in Physics Vol. 18 
	(Springer-Verlag, Berlin, 1993).
%
\bibitem{Carm99} H. J. Carmichael, \textit{Statistical Methods in Quantum 
Optics 1: Master Equations and Fokker-Planck Equations} (Springer-Verlag, 
Berlin, 1999). 
%
\bibitem{Shore} B.~W. Shore, \textit{The Theory of Coherent Atomic 
Excitation, Vol.1} (Wiley, New York, 1990).
%
\bibitem{BeMa11} P.~R. Berman and V.~S. Malinovsky, \textit{Principles 
of Laser Spectroscopy and Quantum Optics} (Princeton University Press, 
Princeton, 2011).
%
\bibitem{CaWa76} H. J. Carmichael and D. F. Walls, J. Phys. B 
	\textbf{9}, 1199 (1976). 
%
\bibitem{Meystre} P. Meystre, \textit{Quantum Optics: Taming The Quantum} 
(Springer-Verlag, Berlin?, 2021). 
%
\bibitem{MeSa4} P. Meystre and M. Sargent III, \textit{Elements of 
Quantum Optics} (Springer-Verlag, Berlin, 2007). 
%
\bibitem{Scully} M. O. Scully and M. S. Zubairy, \textit{Quantum Optics} 
(Cambridge University Press, Cambridge, 1997?).



\end{references}


\end{document}